% Orator ad M. Brutum (orator) --- English translation
% Translated by Claude (Anthropic) from the Latin (Perseus).
% Style guide: STYLE.md.

\ciceroSection{1} Whether it were the harder thing, or the greater, to refuse you when you asked the same of me again and again, or to accomplish what you asked — over this, Brutus, I long and deeply hesitated. For to refuse a man whom I loved beyond all others, and to whom I felt myself most dear, and that when his request was just and his desire honourable — this seemed to me very hard; yet to take up so great a task, one not only difficult to achieve in practice but even to grasp in thought, I scarcely judged to belong to a man who feared the censure of the learned and the wise.

\ciceroSection{2} For what is greater than this: that, when the difference among good orators is so wide, one should judge which is the best species and, as it were, the very shape of speaking? Since you ask this of me again and again, I shall undertake it, not so much in the hope of bringing it off as out of the wish to try; for I would rather, having yielded to your eagerness, be found by you to have fallen short in judgement than, if I had not done it, in goodwill.

\ciceroSection{3} You ask, then — and now more than once — what kind of eloquence I most approve, and what that thing seems to me to be, to which nothing can be added, which I should judge supreme and most perfect. And here I fear that, if I accomplish what you wish and portray that orator whom you seek, I may slow the studies of many, who, weakened by despair, will not wish to try what they distrust their power to attain.

\ciceroSection{4} But it is fitting that all should try everything who have set their hearts on great things, things to be sought with great effort. And if anyone be failed, as it may be, either by his own nature or by that force of surpassing genius, or if he be less furnished with the disciplines of the great arts, let him nonetheless hold to whatever course he can; for to the man who pursues the first place, it is honourable to come to rest in the second or the third. Among poets — to speak of the Greeks — there is room not for Homer alone, or Archilochus, or Sophocles, or Pindar, but also for those second to these, and even for those below the second;

\ciceroSection{5} nor, in philosophy, did the grandeur of Plato deter Aristotle from writing, nor did Aristotle himself, with that admirable knowledge and abundance of his, quench the studies of the rest. Nor were excellent men deterred from the best studies only; not even the craftsmen withdrew from their arts — those who could not match the Ialysus that we saw at Rhodes, nor the beauty of the Coan Venus; nor, daunted by the statue of Olympian Jupiter or by the Spear-Bearer, did the others try the less what they could make, or how far they could advance. So great was the multitude of them, so great in each man's own kind the praise, that, while we marvelled at the highest, we yet approved the lower.

\ciceroSection{6} But among orators — the Greek, at least — it is a marvel how far one man surpasses all the rest; and yet, in the age of Demosthenes, there were many orators great and renowned, and there had been before him, nor did they fail after him. Wherefore there is no reason why the hope of those who have given themselves to the study of eloquence should be broken, or their industry grow faint; for neither is even that which is the best to be despaired of, and among things of high excellence great are those that stand nearest to the best.

\ciceroSection{7} And so, in shaping the supreme orator, I shall fashion such a one as perhaps no man ever was. For I do not ask who there was, but what that thing is, than which nothing can be more excellent — a thing that, in an unbroken stretch of speaking, shines out not often, and I rather think never, yet in some one part shines out at times, more densely in some men, in others perhaps more rarely.

\ciceroSection{8} But I lay it down thus: that there is nothing of any kind so beautiful that there is not something more beautiful still, the source from which it is taken, as an image is drawn, so to speak, from some living face; and this source can be grasped neither by the eyes nor by the ears nor by any sense, but by thought and mind alone. And so of the statues of Phidias, than which we see nothing more perfect of their kind, and of those paintings I have named, we can yet conceive things more beautiful; nor did that craftsman, when he made the form of Jupiter or of Minerva,

\ciceroSection{9} gaze upon some person from whom to draw the likeness, but in his own mind there sat a certain surpassing image of beauty, and, fixing his eyes upon it and intent upon it, he directed his art and his hand to the likeness of it. As, then, in shapes and figures there is something perfect and surpassing, to whose conceived form the things made are referred by imitation, while that form itself falls under no eye, so the form of perfect eloquence we behold with the mind, while its likeness we seek with the ears.

\ciceroSection{10} These forms of things Plato calls \textit{[Greek: ideai]} — that gravest of authorities and teachers not only in conceiving but also in speaking — and he denies that they come into being, and says that they always are and are held together by reason and intelligence; all else is born, perishes, flows, slips away, and does not remain for long in one and the same state. Whatever, then, there is that is disputed by reason and method, that must be brought back to the ultimate form and species of its kind.

\ciceroSection{11} And I see that this first step of mine is drawn not from the disputations of the rhetoricians but fetched from the heart of philosophy, and that, being both ancient and somewhat obscure, it will draw either some reproach or at least some wonder. For men will either marvel what these things have to do with what we are seeking — and these the matter itself, once understood, will satisfy, so that the deep source will not seem fetched without cause — or they will reproach us for tracking down untrodden ways and leaving the beaten paths behind.

\ciceroSection{12} For my part, I am aware that I often seem to be saying new things, when I say things very old but unheard by most; and I confess that I, as an orator — if indeed I am one at all, or whatever I am — came forth not from the workshops of the rhetoricians but from the walks of the Academy. For those are the running-grounds of manifold and varied discourse, in which the footprints of Plato were first impressed. By his disputations and by those of the other philosophers the orator has been both most sharply exercised and most aided; for all the richness, the very forest, so to speak, of speaking, is drawn from them — yet not sufficiently equipped for cases in the Forum, which, as those very men used to say, they left to the ruder Muses.

\ciceroSection{13} Thus this eloquence of the Forum, scorned and rejected by the philosophers, lacked indeed many great aids; but, adorned all the same with words and with thoughts, it won applause among the people and did not fear the judgement or the censure of the few. So the learned wanted popular eloquence, and the eloquent wanted elegant learning.

\ciceroSection{14} Let it be laid down, then, first of all — a thing that will be better understood hereafter — that without philosophy the eloquent man we seek cannot be brought into being; not that all things are in philosophy, but that it gives such help as the training-ground gives the actor; for small things are often most rightly compared with great. For no one can speak with greater breadth and abundance on great and various matters without philosophy —

\ciceroSection{15} since indeed, even in Plato's \textit{Phaedrus}, Socrates says that Pericles excelled the other orators in this, that he had been a hearer of the natural philosopher Anaxagoras; from whom, Socrates holds, besides learning certain other splendid and lofty things, he became rich and fertile, and knew — which is the greatest thing in eloquence — by what modes of speech each part of men's minds is to be moved; and the same may be judged of Demosthenes, from whose letters

\ciceroSection{16} one may gather how constant a hearer of Plato he was — nor indeed, without the discipline of the philosophers, can we discern the genus and species of each thing, nor unfold it by definition, nor distribute it into parts, nor judge what is true and what false, nor discern consequences, see contradictions, distinguish ambiguities. And what shall I say of the nature of things, the knowledge of which furnishes great abundance to oratory — of life, of duties, of virtue, of morals? Can these be spoken or even understood well enough without much study of the very things themselves?

\ciceroSection{17} To all these matters, so many and so great, must be added ornaments beyond number; and these alone, in those days at least, were what the so-called masters of speaking handed down. From which it comes about that no one attains that true and finished eloquence, because the discipline of conceiving is one thing and the discipline of speaking another, and from some men the knowledge of matter, from others the knowledge of words, is sought.

\ciceroSection{18} And so M. Antonius, to whom the age of our fathers gave even the first place in eloquence, a man by nature most penetrating and shrewd, says in that one book which he left behind that he had seen many fluent men, but not one eloquent. Plainly there sat in his mind an image of eloquence, which he discerned with the mind but did not see in actual fact. A man of the keenest genius — for such he was — finding much wanting both in himself and in others, he saw no one who could rightly be called eloquent. And if he thought neither himself nor L. Crassus eloquent,

\ciceroSection{19} then surely he held some form of eloquence grasped in his mind, to which, since nothing was lacking, he could not admit those in whom one thing, or more, was wanting. Let us track him down, then, Brutus, if we can — him whom Antonius never saw, or who in truth never existed at all; and if we cannot imitate and portray him, which that same man said was scarcely granted even to a god, we shall perhaps be able to say what manner of man he ought to be.

\ciceroSection{20} There are, in all, three kinds of speaking, in each of which singly certain men have flourished, but in all of them alike — which is what we want — very few. For there have been the grand-styled, so to speak, with ample and weighty thought and majesty of words, vehement, varied, copious, grave, equipped and ready to move minds and turn them about — a thing which some achieved by a rough, harsh, bristling speech, neither rounded nor brought to a close, others by one smooth, well-built, and finished; and on the other side the slender, the keen, teaching everything and making it clearer rather than fuller, polished by a certain fine and compressed speech; and within that same kind some are shrewd, but unpolished and of set purpose like to the rough and unschooled, while others, in the same leanness, are more harmonious — that is, witty, blooming even, and lightly adorned.

\ciceroSection{21} But there is also a certain middle kind, set between these and, as it were, tempered, using neither the keenness of the latter nor the thunderbolt of the former, neighbour to both, excelling in neither, sharing in each — or rather, if we seek the truth, lacking each. This one flows, as they say, in a single even strain in its speaking, bringing nothing but smoothness and evenness, or it adds a few flowers, as in a garland, and marks out the whole speech with moderate ornaments of words and thoughts.

\ciceroSection{22} Whoever have attained the power of these several kinds, each in its own kind, have won a great name among orators. But the question is whether they have brought off what we want. For we see that there have been some who could speak, in the same person, with ornament and weight, and again with subtlety and finesse. And would that among Latin speakers we could find the likeness of such an orator! It would be a fine thing not to seek what is foreign, but to be content with what is our own.

\ciceroSection{23} But I, the very man who in that conversation of ours set out in the \textit{Brutus} gave much to the Latins — whether to encourage others or because I loved my own countrymen — recall that I set one man, Demosthenes, far before them all; and him I would fit to that eloquence which I conceive in my mind, not to that which I have come to know in any actual man. There has been no one graver than he, none more shrewd, none more measured. And so those must be warned whose untaught talk has grown common — who either wish to be called Attic, or themselves wish to speak in the Attic way — that they should admire this man above all, than whom I do not believe even Athens itself was more Attic; let them learn what the Attic thing is, and measure his eloquence by his strength, not by their own weakness.

\ciceroSection{24} For at present each man praises only so much as he hopes himself to be able to imitate. Yet I think it not amiss to teach those who are endowed with the best zeal but a less firm judgement what the proper praise of the Attic writers is. The eloquence of orators has always been governed by the discernment of its hearers. For all who wish to win approval look to the will of those who hear, and to their will, their pleasure, their nod, they fashion and accommodate themselves wholly.

\ciceroSection{25} And so Caria and Phrygia and Mysia, being the least polished and the least elegant, took to themselves, as suited to their ears, a certain rich and, so to speak, fatted kind of speech; which their neighbours the Rhodians, with no wide sea set between, never approved, and Greece much less, while the Athenians rejected it utterly. Their judgement was always discerning and pure, so that they could endure to hear nothing but what was uncorrupted and elegant. And while the orator was a servant to their scruple, he dared to set down no word unusual, no word distasteful.

\ciceroSection{26} And so that man whom we have said surpassed the rest, in that speech \textit{For Ctesiphon} — by far his best — at the outset spoke more quietly; then, while arguing about the laws, more compressedly; afterward, gradually kindling the jurors, as soon as he saw them ablaze, in the rest of the speech he leapt out more boldly. And yet, even with him weighing carefully the weight of every word, Aeschines reproaches certain expressions and assails them, and mockingly calls them harsh, distasteful, unbearable; nay, he even asks the man himself — when indeed he calls him a brute — whether those are words or monstrosities; so that to Aeschines not even Demosthenes seems to speak in the Attic way.

\ciceroSection{27} For it is easy to mark some word that blazes, so to speak, and then, once the fires of the mind are quenched, to laugh at it. And so, clearing himself, Demosthenes makes a jest of it: he denies that the fortunes of Greece were staked on this — whether he had used this word or that, whether he had stretched out his hand this way or that. How, then, would a Mysian or a Phrygian be heard at Athens, when even Demosthenes is assailed as overblown? But when one had begun to chant in the Asiatic manner, with a falling, howling voice, who would bear him — or rather, who would not order him carried off?

\ciceroSection{28} Those, then, who accommodate themselves to the polished and scrupulous ears of the Attic writers are to be reckoned to speak in the Attic way. Of these there are several kinds; but these men suspect only the one, of whatever sort it is. For they think that he who speaks roughly and without cultivation, provided he do it elegantly and clearly, he alone speaks in the Attic way. They are wrong in the "alone"; in the "Attic" they are not deceived.

\ciceroSection{29} For by their judgement, if that alone is Attic, then not even Pericles spoke in the Attic way, to whom the first place was awarded without dispute; and had he used the slender kind, he would never have been said by the poet Aristophanes to flash, to thunder, to throw all Greece into confusion. Let that most charming and most polished writer Lysias, then, speak in the Attic way — for who could deny it? — provided we understand that the Attic thing in Lysias is not that he is slender and unadorned, but that he has nothing unusual or inept; while to speak with ornament, with weight, and with abundance is either the mark of the Attic writers, or else let neither Aeschines nor Demosthenes be Attic.

\ciceroSection{30} But behold, certain men profess themselves to be Thucydideans: a new kind of the unschooled, and unheard-of. For those who follow Lysias follow a kind of pleader — not, to be sure, ample and grand, yet subtle and elegant, and one who could make a fine stand in cases in the Forum. But Thucydides narrates deeds done, and wars, and battles, gravely indeed and well, but nothing can be carried over from him to the use of the Forum and of public affairs. Those very harangues of his have so many obscure and hidden thoughts that they can scarcely be understood — which in civil oratory is the very greatest fault of all.

\ciceroSection{31} But what perversity is there in men so great that, with the harvest discovered, they should still feed on acorns? Could the diet of mankind be refined by the kindness of the Athenians, and could speech not be refined too? And then — what Greek rhetorician ever drew anything from Thucydides? ``But he was praised by everyone.'' I grant it; yet praised as a wise, austere, and weighty expositor of events — not as a man who handled cases in the courts, but as one who narrated wars in his histories. And so he was never reckoned an orator.

\ciceroSection{32} Nor indeed, had he not written history, would his name survive — and that although he had held office and was of noble birth. Yet no one imitates the weight either of his words or of his thoughts; rather, when men have uttered certain maimed and gaping phrases, which they could have managed even without a teacher, they fancy themselves true-born Thucydideses. I have met, too, with one who longed to be like Xenophon, whose speech is indeed sweeter than honey, but as far as can be from the din of the Forum.

\ciceroSection{33} Let us turn back, then, to the man we wish to begin and to mould at last with that eloquence which Antonius found in no one. We are attempting a great work, Brutus, and a hard one; but to one who loves, I think nothing is difficult. And I love, and have always loved, your talent, your studies, your character. Indeed I am kindled more each day — not by longing alone, though that wears me away, as I miss our meetings, the companionship of daily life, your most learned conversation — but also by the incredible report of your admirable virtues, which, unlike in appearance, are joined together by wisdom.

\ciceroSection{34} For what is so distant from austerity as affability? Yet who was ever held either more upright than you or more agreeable? What is so difficult as to be loved by all when adjudicating the disputes of the many? Yet you bring it about that you send away even those against whom you decide reconciled and at peace. And so you contrive that, although you do nothing for the sake of favour, all that you do is nonetheless welcome. Therefore, of all lands, Gaul alone does not blaze with the common conflagration; there you enjoy your own company, while you are recognized in the full light of Italy and move among the very flower and strength of our best citizens. And how great a thing is this — that amid the greatest occupations you never interrupt the studies of learning, but are always either writing something yourself or summoning me to write!

\ciceroSection{35} And so I undertook this task at once, as soon as the \textit{Cato} was finished — a work I would never have touched, fearing times hostile to virtue, had I not judged it an offence not to obey you when you urged me and stirred up in me the memory of that man so dear to me — but I call you to witness that it is at your asking, and against my own refusals, that I have dared to write these things. For I wish the charge to be shared between us, so that, if I should prove unable to sustain so great an inquiry, the fault of an unjust burden laid on me may be yours, and mine the fault of taking it up; though even so the praise of the gift you bestowed will make good any error in my judgment.

\ciceroSection{36} But in every matter it is most difficult to set out the form — what the Greeks call \textit{[Greek: charakter]} — of the best, because one thing seems best to one man, another to another. ``I delight in Ennius,'' says someone, ``because he does not depart from the common usage of words''; ``I in Pacuvius,'' says another: ``in him all the verses are ornate and elaborated, in the other far more carelessly done''; suppose yet another prefers Accius. For judgments vary, as among the Greeks, nor is it easy to explain which form most excels. In paintings some are delighted by the rough, the uncultivated, the recessed and shadowed; others, on the contrary, by the bright, the cheerful, the well-lit. What is there from which you might draw out some prescription or formula, when each thing excels in its own kind, and the kinds are many? By this scruple I was not driven back from the attempt; I judged that in all things there is something best, even if it lay hidden, and that it could be judged by one who was expert in the matter.

\ciceroSection{37} But since there are several kinds of oratory, and these diverse, and they do not all fall under a single form — those of panegyrics and of histories and of such persuasive pieces as the \textit{Panegyricus} which Isocrates composed, and many others of those who are called sophists, and the form of the remaining kinds of composition, which stand apart from forensic contention, and of that whole class which the Greeks name \textit{[Greek: epideiktikon]}, because it is got up, as it were, for inspection, for the sake of giving pleasure — this I shall not embrace at present; not that it is to be neglected, for it is, as it were, the nurse of that orator whom we wish to mould and of whom we are labouring to say something more exacting. From this both his fullness of words is fed, and the arrangement of them, and his rhythm enjoys a freer kind of license.

\ciceroSection{38} Indulgence, too, is granted to the trim balancing of thoughts, and pointed, exact, and rounded periods of words are conceded; and the work is done deliberately — not from ambush, but openly and in plain sight — so that words may answer to words as if measured out and matched, so that opposites may often be set against each other and contraries compared, and so that clauses may end alike and bring back the same sound in their cadence. These things, in the truth of real cases, we do both far more rarely and certainly more covertly. In the \textit{Panathenaicus}, however, Isocrates confesses that he pursued these effects studiously; for he had written not for the contest of the courts, but for the pleasure of the ears.

\ciceroSection{39} They say that Thrasymachus of Calchedon first handled these matters, and Gorgias of Leontini, and after them Theodorus of Byzantium and many others whom Socrates, in the \textit{Phaedrus}, calls word-artificers \textit{[Greek: logodaidalous]}; of whose work much is pointed enough, but, as being just then first coming into being, minute and like little verses, and overly painted. The more wonderful, by contrast, are Herodotus and Thucydides; for although their age had fallen within the times of those I have named, they themselves kept very far from such dainties — or rather such follies. The one flows, without any rough places, like a calm river; the other is borne along more impetuously, and of warlike matters sings even, in a manner, a war-song. And by these first, as Theophrastus says, history was stirred to dare to speak more richly and more ornately than its predecessors.

\ciceroSection{40} To the age of these men succeeded Isocrates, who beyond all others of the same kind is always praised by us — though now and then, Brutus, you gently and learnedly resist. But perhaps you will grant me the point, if you have learned what it is in him that I praise. For when Thrasymachus seemed to him chopped up with minute rhythms, and Gorgias likewise — though these two are reported to have been the first to join words by a certain art — and Theodorus, again, too abrupt and not, so to speak, sufficiently rounded, he was the first to set about expanding his thoughts with words and filling them out with softer rhythms; and since in this he taught those who became foremost partly in speaking and partly in writing, his house was held to be the workshop of eloquence.

\ciceroSection{41} And so, just as I, when I was praised by our Cato, easily bore to be censured by the rest, so Isocrates may rightly, on the testimony of Plato, despise the judgments of others. For there is, as you know, near the very end of the \textit{Phaedrus}, Socrates speaking in these very words: ``Isocrates is still a young man, Phaedrus; but I should like to say what I forecast of him.'' ``What, then?'' said the other. ``He seems to me of too great a talent to be compared with the speeches of Lysias; and besides, of a greater natural bent toward virtue; so that it will be no wonder at all if, when he has advanced in years, he should either, in this kind of oratory to which he now devotes himself, surpass all who have ever touched oratory by as much as he now surpasses other boys, or, if he is not content with these things, should reach by some divine impulse of mind toward greater ones; for there is by nature in this man's mind a certain philosophy.''

\ciceroSection{42} Such are the things Socrates forecasts of him as a young man. But Plato writes these things of him when he was older — and he writes as his contemporary, and indeed, a railer at all rhetoricians, he admires this one alone; so let those who do not love Isocrates suffer me to err along with Socrates and with Plato. The kind of oratory, then, is sweet, and loose and flowing, pointed in its thoughts and resounding in its words, in that epideictic class which we have said is proper to the sophists — fitter for parade than for battle, dedicated to the gymnasium and the wrestling-school, scorned and driven out of the Forum. But since eloquence, reared on this nourishment, afterward gives itself its own colour and strength, it was not out of place to speak of the orator's cradle, as it were. But these things belong to the games and the parade; let us, for our part, now come to the line of battle and the fight.

\ciceroSection{43} Since three things are to be looked to by the orator — what he is to say, and in what place each thing is to be said, and in what manner — we must indeed say what is best in each; but somewhat otherwise than is usual in the handing down of an art. We shall lay down no precepts, for that is not what we have undertaken, but we shall sketch the appearance and form of surpassing eloquence; nor shall we set out by what means it is procured, but only of what sort it seems to us to be.

\ciceroSection{44} And the first two briefly; for they are not so much marks of the highest distinction as they are necessary, and yet they are almost shared with many. For both to discover and to judge what you are to say are indeed great matters, and are, as it were, the soul within the body, but they belong more to good sense than to eloquence; and yet in what case is good sense without a place? Let this orator, then — the one we wish to be supreme — know the topics of arguments and of reasonings.

\ciceroSection{45} For since, whatever it is that is engaged in controversy or in contention, in it the question is either whether a thing is, or what it is, or of what quality it is — whether it is, by signs; what it is, by definitions; of what quality it is, by the parts of right and wrong — that the orator may be able to use these (not that common orator, but this surpassing one), let him always, if he can, call the controversy away from particular persons and times; for it is permitted to debate more broadly of the class than of the part, so that what has been proved in the whole must necessarily be proved in the part —

\ciceroSection{46} this question, then, transferred from particular persons and times to the account of the whole class, is called a thesis \textit{[Greek: thesis]}. In this Aristotle trained his young men, not in the manner of philosophers, to discourse thinly, but to the fullness of the rhetoricians, on both sides, so that the thing might be said more ornately and more richly; and the same man handed down the topics — for so he calls them — as, so to speak, the marks of arguments, from which all oratory might be drawn on either side.

\ciceroSection{47} This orator of ours, then, will so act — for we are seeking not some declaimer from the school or a brawler from the Forum, but the most learned and most perfect of men — that, since fixed topics are handed down, he may run through them all, use the apt ones, and speak by classes; from which there flow also what are called the common topics. Nor indeed will he use this abundance imprudently, but will weigh everything and pick and choose; for the weight of the same arguments, drawn from the same topics, is not always nor in all cases the same.

\ciceroSection{48} He will therefore bring judgment to bear, and will not only discover what to say but will also weigh it. For nothing is more fertile than men's talents, especially those that have been cultivated by the disciplines. But just as fruitful and abundant fields pour forth not only the crops but also weeds most hostile to the crops, so sometimes from those topics there are born things that are either trivial, or foreign to the case, or of no use. Unless a great discernment of these is brought to bear by the orator's judgment,

\ciceroSection{49} how shall he cling to his strong points and dwell among them — or soften the harsh ones, or hide what cannot be explained away, and altogether suppress it, if he is allowed, or lead the minds elsewhere, or bring up some other thing which, set against it, may be more plausible than that which will stand in his way?

\ciceroSection{50} And then, with what care will he arrange the things he has discovered? For that was the second of the three. He will surely make the porches of his case honourable and its approaches striking; and when he has seized the minds at the first onset, he will weaken and shut out what is contrary; of his strongest points he will set some first, some last, and will hammer home the slighter ones in between.

\ciceroSection{51} And in the first two parts of speaking we have described, summarily and briefly, of what sort he would be. But, as was said before, in these parts, though they are weighty and great, there is less of art and of labour; whereas, when he has discovered both what to say and in what place, that is by far the greatest thing — to see in what manner. For it is a shrewd saying, which our Carneades used to make, that Clitomachus said the same things, but Charmadas said them in the same manner too. And if in philosophy it matters so much how you speak — where the substance is regarded and the words are not weighed — what, after all, must we judge in cases, which are governed entirely by the speaking?

\ciceroSection{52} And this indeed, Brutus, I gathered from your letter: that you were not inquiring of me what kind of orator I should wish to be supreme in discovery and in arrangement, but you seemed to me to be asking this — which kind of speech itself I judged the best: a difficult matter, by the immortal gods, and the most difficult of all. For since speech is soft and tender and so pliant that it follows wherever you turn it, the varied natures and dispositions of men have produced kinds of speaking that differ widely among themselves.

\ciceroSection{53} A stream of words and a fluency are dear to the hearts of some, who set eloquence in the speed of speech; others are delighted by intervals distinct and punctuated, by pauses and breathings: what could be so different? And yet in each there is something that excels. Some labour at smoothness and evenness and a kind of pure, white style of speaking; and behold, others pursue a certain hardness and severity in their words and, as it were, a mournfulness in their speech; and, as we divided a little before — that some wished to seem grand, others plain, others measured — there are found exactly as many kinds of orators as we said there are kinds of oratory.

\ciceroSection{54} And since I have now begun to swell this office more amply than was demanded of you — for when you asked only about the kind of speech, I answered, even if briefly, about discovery and arrangement too — let me now speak not only about the manner of speech but also about the manner of delivery: so that no part will be passed over, since on the subject of memory there is nothing to be said in this place, that being common to many arts.

\ciceroSection{55} Now the manner in which a thing is said lies in two things — in the delivery and in the diction. For delivery is, as it were, a kind of eloquence of the body, since it consists of voice and movement. The variations of the voice are exactly as many as those of the soul, which is most of all moved by the voice. And so that perfect orator, whom this discourse of ours has long been pointing to, will, however he wishes himself to seem affected and the mind of the hearer to be moved, bring to bear a fixed tone of voice to that end; about which I would say more, if this were the time for laying down precepts, or if you were asking it. I would speak too of gesture, with which the face is joined; in all of which it can scarcely be said how much it matters in what manner the orator uses them.

\ciceroSection{56} For both the inarticulate have often gathered the fruit of eloquence by the dignity of their delivery, and many eloquent men have been thought inarticulate through ugliness of delivery; so that it was not without cause that Demosthenes assigned to delivery the first place, and the second, and the third. For if there is no eloquence without it, while it, without eloquence, is of such power, then surely it counts for very much in speaking. He, then, who would seek the first place in eloquence will wish both to speak fiercely with a strained voice and gently with a lowered one, and to seem grave when his voice sinks and pitiable when it bends;

\ciceroSection{57} for there is a wonderful nature in the voice, of which, out of three tones in all — the bending, the high, and the low — so great and so sweet a variety has been perfected in song. There is, moreover, even in speaking a certain more hidden melody — not this epilogue of the rhetoricians from Phrygia and Caria, which is almost a chant, but that which Demosthenes and Aeschines mean, when the one casts at the other his inflections of voice; Demosthenes even says that the man wailed whom he repeatedly says had a voice sweet and clear.

\ciceroSection{58} On which point this too seems to me worth noting, toward the pursuit of sweetness in the voice: that nature herself, as though she were tuning the speech of men, has placed in every word the acute tone, and not more than one, nor farther back from the last syllable than the third; so that industry should the more follow nature as its guide to the pleasure of the ears.

\ciceroSection{59} And the goodness of the voice, indeed, is to be wished for; for it is not in our power, but the management and use of it is in our power. Therefore that foremost orator will vary and change it: he will run through all the degrees of sound, now straining, now relaxing them. And he will so use movement that nothing is superfluous. In gesture, the posture upright and lofty; the pacing rare and not too long; the advance moderate and that rarely; no softness of the neck, no nimbleness of the fingers, no joint falling to a beat; ruling himself rather by his whole trunk and by a manly bending of the flanks, by a thrusting forth of the arm in heated passages and a drawing in of it in the calmer ones.

\ciceroSection{60} As for the face, which after the voice has the most power, what dignity, and what charm besides, it will bring! And when in it you have brought it about that there is nothing inept or grimacing, then there is a certain great control of the eyes. For as the face is the image of the soul, so the eyes are its interpreters; and the measure of their cheerfulness and, in turn, of their sadness will be tempered by the very matters that are being treated.

\ciceroSection{61} But now I must sketch the likeness of that perfect orator and of supreme eloquence. That he excels in this one thing alone — in speech — while everything else lies hidden within him, the very name declares: for the man who has embraced all these arts is not called inventor or arranger or performer, but in Greek, from the act of speaking out, \textit{[Greek: rhêtôr]}, and in Latin "eloquent." Of the other endowments that belong to the orator each man claims some share for himself; but the highest power of speaking — that is, of expressing in words — is granted to him alone.

\ciceroSection{62} For although certain philosophers too have spoken with ornament — since Theophrastus drew his very name from the divinity of his speech, and Aristotle challenged Isocrates himself, and they say the Muses spoke as if with the voice of Xenophon, and Plato stood out by far above all who have ever written or spoken, first in both sweetness and weight — still, the speech of these men has neither the sinews nor the stings of the orator and of the courts.

\ciceroSection{63} They talk with the learned, whose minds they prefer to soothe rather than to rouse; and they talk of subjects untroubled and quite unstormy, for the sake of teaching, not of capturing — so that in the very thing whereby they go fishing for some pleasure in their speaking, they may seem to some to do more than is needful. From this kind, then, it is not hard to distinguish the eloquence with which we are now concerned.

\ciceroSection{64} For the speech of the philosophers is soft and shadowed, furnished with neither thoughts nor words drawn from the people, nor bound by rhythms, but loosed more freely; it has nothing angry, nothing spiteful, nothing fierce, nothing piteous, nothing cunning — chaste, modest, an unspoiled maiden, so to speak. And so it is called conversation rather than oratory. For although all speech is "speaking," still the speaking of the orator alone is marked by this special name.

\ciceroSection{65} The likeness of the Sophists, of whom I spoke above, seems to need sharper distinguishing, since they all want to pursue those same flowers that the orator employs in his cases. But they differ in this: that, since their aim is not to trouble men's minds but rather to calm them, and not so much to persuade as to delight, they do this both more openly than we do and more often; they seek out thoughts that are neatly turned rather than convincing, they often wander from the matter, they weave in tales, they transfer words too far afield and arrange them as painters arrange a variety of colours, they match equal with equal, set opposite against contrary, and very often round off the ends of clauses alike.

\ciceroSection{66} Close to this kind is history, in which the narrative too is given with ornament, and often a region or a battle is described; speeches and exhortations are even inserted, but in these a certain drawn-out and flowing style is sought, not this twisted and keen oratory. From these writers this eloquence which we seek is to be called away, not much otherwise than from the poets. For the poets too have raised the question what it was that set them apart from the orators: rhythm and verse seemed chiefly to do so in earlier times, but now among the orators rhythm itself has grown common.

\ciceroSection{67} For whatever falls under some measure of the ear — even if it is wanting in verse, for that, in oratory, is indeed a fault — is called rhythm, which in Greek is termed \textit{[Greek: rhythmos]}. And so I see it has seemed to some that the speech of Plato and Democritus, though it is wanting in verse, is yet to be reckoned poetry rather than that of the comic poets, because it moves with more impetus and uses the most brilliant lights of words; whereas in the comic poets there is nothing, save that they have their little verses, that differs from everyday conversation. And yet this is not what is greatest in a poet — although he is the more praiseworthy for pursuing the orator's virtues even while he is more closely bound by verse.

\ciceroSection{68} For my own part, even if the voice of certain poets is grand and ornamented, still I hold that in it there is a greater licence than in us for making and joining words; and also that not a few of them serve their wish more by words than by matter. Nor indeed, if there is some one thing alike between them — and that is the judgment and choice of words — can the unlikeness of everything else for that reason not be perceived. But this is beyond doubt; and if it holds any matter for question, that very point is yet not necessary to what is set before us. The orator, then, set apart from the eloquence of the philosophers, of the Sophists, of the historians, of the poets, must be expounded by us as to what he is to be.

\ciceroSection{69} He will be eloquent, then — for it is this man we seek, on Antonius' authority — who in the Forum and in civil cases will so speak as to prove, to delight, to sway. To prove is a matter of necessity, to delight of charm, to sway of victory: for this one thing of all has the most power toward winning cases. But as many as are the duties of the orator, so many are the kinds of speaking: the plain in proving, the moderate in delighting, the vehement in swaying; and in this one thing lies all the orator's force.

\ciceroSection{70} He, then, who is to govern and, as it were, temper this threefold variety will need great judgment and the highest faculty besides; for he will both decide what each thing requires and be able to speak in whatever way the case shall demand. But the foundation of eloquence, as of all else, is wisdom. For as in life, so in speech, nothing is harder than to see what is fitting. The Greeks call this \textit{[Greek: prepon]}; let us by all means call it "propriety." On it many fine precepts are given, and the matter is most worthy of study; through ignorance of it men go astray not only in life but very often in poetry and in oratory too.

\ciceroSection{71} Now what is fitting must be looked to by the orator not only in his thoughts but also in his words. For not every fortune, not every rank, not every authority, not every age — nor indeed every place or season or hearer — is to be handled with the same kind of words or of thoughts; and always, in every part of a speech as of life, one must consider what is fitting. This depends both on the matter under discussion and on the persons both of those who speak and of those who hear.

\ciceroSection{72} And so the philosophers are wont to treat this topic — which extends far and wide — in their discourses on duties, though not when they dispute about the right itself, for that is a single thing; the grammarians treat it in the poets; the eloquent in every kind and every part of cases. For how unfitting it is, when you are pleading about rain-water dripping from the eaves before a single judge, to use the most ample words and commonplaces, or to speak of the majesty of the Roman people in a lowered and plain manner! Here men err in the whole kind, but others err in the person — their own, or the judges', or even their adversaries' — and not only in matter but often in a word; for although without matter a word has no force, still the same matter is often either approved or rejected as it is set forth in one word and another.

\ciceroSection{73} And in all things one must see how far to go; for although each thing has its own measure, still too much offends more than too little. In this Apelles used to say that those painters too went wrong who did not feel what was enough. This is a great topic, Brutus — as is not lost on you — and it calls for another great volume; but for what is at hand, this much is enough. When we say that this is fitting — a thing we constantly invoke in all that is said and done, in matters least and greatest — when, I say, we say that this is fitting and that is not fitting, the magnitude of the principle appears at every turn, and it is one thing and a wholly other whether you say "fitting" or "obligatory." For "obligatory" declares the perfection of duty,

\ciceroSection{74} which is always to be observed and by all; whereas "fitting" is, as it were, to be apt and consonant with the season and the person. This holds good very often in deeds and likewise in words, and finally in the face and gesture and gait — and "unfitting," conversely, in the same. If the poet flees this as the greatest of faults — the poet who goes wrong even when he frames the speech of an upright man for a base one, or of a wise man for a fool; if, finally, that painter saw, when at the sacrifice of Iphigenia Calchas was sorrowful, Ulysses more sorrowful, Menelaus grieving, that Agamemnon's head must be veiled, since he could not represent that supreme grief with his brush; if, finally, the actor seeks out what is fitting — what are we to think the orator must do? But since this is so great a thing, let the orator see for himself what he is to do in his cases and in their members, so to speak: this at least is clear, that not only the parts of a speech but whole cases too are to be handled, one in one form of speaking, another in another.

\ciceroSection{75} It follows that the mark and pattern of each kind should be sought: a great task and a hard one, as I have already often said; but for those setting out, what we were undertaking had to be weighed, while now indeed, wherever we are carried, the sails must surely be spread. And first the one whom alone, in truth, they call "Attic" must be formed by us.

\ciceroSection{76} He is lowered and unassuming, imitating ordinary usage, differing from the unskilled more in fact than in seeming. And so those who hear him, however tongue-tied they themselves may be, are nonetheless confident that they can speak in that manner. For the plainness of such oratory does indeed seem imitable to one who judges, but nothing is less so to one who tries. For although it is not of the fullest blood, still it ought to have some sap, so that, even if it lacks those greatest powers, it may be, so to speak, of sound health.

\ciceroSection{77} First, then, let us free him as if from the bonds of rhythm. For there are, as you know, certain oratorical rhythms, of which we shall treat presently, to be observed by a certain method — but in another kind of speaking; in this one they are to be wholly relinquished. Let there be something loosed, yet not adrift, so that he may seem to walk freely, not to wander licentiously. Let him neglect even to cement word to word, as it were. For that gaping and clashing of vowels, so to speak, has a certain softness in it, and something that betokens a not unwelcome carelessness in a man who labours over his matter more than over his words.

\ciceroSection{78} But there will be the rest to look to, once these two have been left more free to him: the periodic structure and the gluing-together of words. For those very contracted and minute things are not to be handled carelessly, but there is even a kind of carelessness that is careful. For as some women are said to be unadorned, as though that very thing became them, so this plain oratory pleases even when uncombed; for in both there happens something whereby it is the more charming, but not so as to show. Then all conspicuous ornament, like that of pearls, will be removed;

\ciceroSection{79} not even the curling-irons will be applied; all the cosmetics of a false fairness and flush will be repelled; only elegance and neatness will remain. The diction will be pure and Latin, the speech lucid and plain, and what is fitting will be looked round for; one thing will be absent — that which Theophrastus reckons fourth among the merits of speech: ornament, that sweet and abundant thing. Sharp and frequent thoughts will be set down, and dug out from somewhere in concealment; and — what will rule in this orator — there will be a modest use of the orator's furniture, so to speak.

\ciceroSection{80} For there is in a manner our own furniture, which lies in ornaments — one kind belonging to matters, another to words. The ornament of words is twofold: one of single words, the other of words in combination. The single word is approved among words proper and customary when it either sounds best or makes the matter clearest; among words not proper, when it is transferred and made from elsewhere, as if on loan, or made by the speaker himself and new, or ancient and unaccustomed; but even the unaccustomed and ancient are among the proper, save that we use them rarely.

\ciceroSection{81} Words in combination, on the other hand, have ornament if they produce some symmetry — something that does not survive a change of words while the thought survives. For the ornaments of thought, which abide even if you have changed the words, are indeed very many, but those that stand out are fewer. Therefore that slender orator, provided he be elegant, will neither be bold in coining words and will be modest and sparing in transferring them and in using ancient ones, and rather restrained in the rest of the ornaments both of words and of thoughts; in metaphor he will perhaps be the more frequent, since all speech uses it most often, not only of city folk but even of country folk — seeing that it is theirs to say that the vines "bud with gems," that the fields "thirst," that the crops "are glad," that the grain "runs riot."

\ciceroSection{82} None of these is too bold, since the term is either like that from which you transfer it, or, if the thing has no name of its own, is taken for the sake of teaching, not of play. Of this ornament our plain man will make a little freer use than of the rest, yet not so licentiously as if he were using the most ample kind of speaking; and so the unfitting — whose nature must be understood from the fitting — appears here too, when some word is transferred too loftily and is set in lowly speech, the very word that in another style would be fitting.

\ciceroSection{83} As for that symmetry which lights up the arrangement of words with those lights which the Greeks call the gestures of speech, as it were — \textit{[Greek: schêmata]}, the same word being transferred by them also to the ornaments of thought — this plain orator does indeed employ it (whom some, however wrongly in other respects, rightly call "Attic" on this one ground alone), but a little more sparingly; for, just as at a banquet's furnishing one who draws back from magnificence will wish to seem not only sparing but also elegant, so he will choose what he is to use.

\ciceroSection{84} For most things suit the very thrift of this orator of whom I speak. For those things of which I spoke before are to be shunned by this keen man: clauses matched with clauses and closed off alike and falling in the same fashion, and charms got, as it were, by the changing of a letter — lest a laboured symmetry and a manifest hunting after delight be caught in the act and show.

\ciceroSection{85} Likewise, if any repetitions of words call for some straining of the voice and an outcry, they will be alien to this lowered manner of speaking; the rest he will be able to use without restriction, only let him loosen and divide the continuous run of words, and use words as customary as may be, and metaphors as soft as may be; he will take up even those lights of thought which will not be vehemently brilliant. He will not bring the commonwealth on to speak, nor raise the dead from the underworld, nor heap up many points crowded together and bind them in a single embrace. These belong to stronger lungs, and are neither to be expected nor demanded of the man we are forming; for he will be, in his speech as in his voice, the more subdued.

\ciceroSection{86} But most of the things from those grander styles will suit even this slenderness, although he will use the same ornaments more roughly; for such is the man we bring forward. His delivery will be added — not tragic nor of the stage, but with a moderate movement of the body, while accomplishing much by the face nonetheless; not in that way by which they say one "pulls a face," but in the one whereby they show with candour in what sense they pronounce each thing.

\ciceroSection{87} On this kind of oratory wit too will be sprinkled, which counts for ever so much in speaking. Of it there are two kinds: one of pleasantry, the other of raillery. He will use both: the one in telling something with charm, the other in casting and launching a jest, of which the kinds are several; but at present we are about another matter.

\ciceroSection{88} This much, however, we do urge: that the orator will so use the ridiculous as to make it neither too frequent, lest it be buffoonish, nor a little obscene, lest it be like the mime, nor petulant, lest it be base, nor aimed at calamity, lest it be inhuman, nor aimed at crime, lest laughter usurp the place of hatred, and as nothing alien either to his own person or to the judges' or to the occasion. For these are all referred to that unfitting we spoke of.

\ciceroSection{89} He will also shun jests that are sought out and fabricated not on the spur of the moment but brought from home, which for the most part fall flat. He will spare both friendships and high stations, will shun incurable insults, will only strike at his adversaries — and not even those always, nor all of them, nor in every way. With these exceptions he will so use wit and pleasantry that, for my part, I have known no such man among those new Atticists of yours, though that is surely Attic above all else.

\ciceroSection{90} This I judge to be the form of the plain orator — yet a great one too, and a genuine Attic; since whatever is salty or wholesome in a speech is the property of the Attics. Yet of them not all are witty: Lysias is sufficiently so, and Hyperides; Demades is reported beyond the rest; Demosthenes is held to have less of it — though to me nothing seems more urbane than he, but he was not so given to raillery as to pleasantry. The one is the mark of a keener wit, the other of a greater art.

\ciceroSection{91} The second style is fuller and considerably more robust than the plain one of which I have spoken, yet more restrained than the grandest, of which I shall speak presently. In this kind there is the least, perhaps, of sinew, but the most of sweetness. For it is richer than the lean style, and more restrained than the ornate and copious one.

\ciceroSection{92} To this style all the ornaments of speech are suited, and in this form of oratory there is the greatest charm. In it many flourished among the Greeks, but in my judgment Demetrius of Phalerum surpassed the rest; for while his speech flows on calmly and quietly, it is at the same time lit up, as if by certain stars, with words transferred and altered. By transferred I mean — as I have often said by now — those words which by likeness are carried over from another thing, whether for the sake of sweetness or out of poverty of language; by altered, those in which, in place of the proper word, another is substituted which means the same, drawn from some kindred thing.

\ciceroSection{93} And although this too is done by transference, yet Ennius transferred in one way when he said \textit{I am bereft of citadel and city}, and would have transferred in another had he said ``citadel'' for ``country''; and when he says \textit{rugged Africa trembles with a fearful tumult}, he uses ``Africa'' for ``the Africans'' by alteration. This figure the rhetoricians call \textit{hypallage} \textit{[Greek: hypallagē]}, because words are as it were exchanged for words; the grammarians call it \textit{metonymy} \textit{[Greek: metōnymia]},

\ciceroSection{94} because nouns are transferred. Aristotle, however, ranks both these under transference, along with what he calls \textit{catachresis} \textit{[Greek: katachrēsis]} — abuse of a word — as when we call a soul ``shrunken'' instead of ``small''; and we misuse neighboring words, if there is need, either because it pleases or because it is fitting. And when several transferences have flowed on without a break, the speech becomes plainly something else; this kind, accordingly, the Greeks call \textit{allegory} \textit{[Greek: allēgoria]} — rightly so in name, but in its genus better understood by the man who calls all such things transferences. Demetrius uses these very frequently, and they are most delightful; and although transference is abundant in him, alterations are nowhere thicker.

\ciceroSection{95} Into this same style of speech — for I am speaking of the moderate and tempered kind — fall all the brilliances of words and many of thought as well; by it broad and learned discussions will be unfolded, and the commonplaces will be delivered without strain. In short, such men generally come out of the schools of the philosophers; and unless that stronger speaker is set up beside him for comparison, this one of whom I speak will win approval on his own.

\ciceroSection{96} For there is besides a certain distinguished and flowering style of speech, painted and polished, in which all the graces of words, all the graces of thought, are woven in. This whole style flowed down into the Forum from the springs of the sophists, but, scorned by the subtle and rebuffed by the weighty, it has settled into that middle ground of which I am speaking.

\ciceroSection{97} The third style is the ample, the copious, the weighty, the ornate, in which surely lies the greatest power. This is the eloquence whose ornament and abundance of speech the admiring nations have allowed to prevail most of all in their states — but this eloquence, the kind that is carried along with a great rushing and a sounding, the kind all looked up to, marveled at, and despaired of attaining for themselves. It belongs to this eloquence to work upon men's minds, to move them by every means. Now it shatters, now it steals into the senses; it implants new convictions, and tears out those that are rooted.

\ciceroSection{98} But there is a great difference between this style of speech and the earlier ones. The man who has labored at that subtle and sharp style, so as to speak shrewdly and pointedly, and has aimed at nothing higher, is, once he has perfected this one thing, a great orator — and, if not the greatest, at least he will tread on the least slippery ground, and, once he has stood firm, will never fall. As for the middle speaker, the one I call moderate and tempered, provided only that he has equipped himself well in that style, he will not dread the doubtful and uncertain hazards of speaking; even if at times he succeeds less, as often happens, he will yet run no great risk — for he cannot fall from any height.

\ciceroSection{99} But this orator of ours, whom we set in the first place — weighty, keen, aflame — if he was born for this one thing alone, or has trained himself in this alone, or has devoted himself to this single kind, and has not tempered his abundance with those other two styles, deserves the utmost contempt. For the restrained speaker, because he speaks shrewdly and with old craft, already passes for wise; the middle one is charming; but this most copious speaker, if he is nothing else, is apt to seem scarcely sane. For a man who can say nothing calmly, nothing gently, nothing in an ordered, defined, distinct, and witty manner — especially when cases must be handled partly whole in that manner, partly in some one part — if such a man, before his hearers' ears are prepared, begins to set the matter ablaze, he seems to rave among the sane and, like a drunken man, to reel as if among the sober.

\ciceroSection{100} We have him, then, Brutus, the man we are seeking — but in our minds; for had I laid hold of him with my hand, not even he, for all his great eloquence, would have persuaded me to let him go. But the eloquent man has surely been found, the one Antonius never saw. Who, then, is he? I shall embrace it briefly, and discuss it at greater length. He is eloquent who can speak of small matters with subtlety, of great matters with weight, and of middling matters with measure. ``No such man,'' you will say, ``ever existed.''

\ciceroSection{101} Suppose he never did. For I am arguing not about what I have seen but about what I desire; and I return to that form and ideal of Plato's, of which I had spoken, which, though we do not perceive it with the eye, we can yet hold in the mind. For I am not seeking an eloquent man, nor anything mortal and perishable, but that very thing itself, whoever possesses which is eloquent — which is nothing other than eloquence itself, which we can see with no eyes but those of the mind. He, then, will be eloquent — to repeat once more that same definition — who can speak of small things in a restrained way, of moderate things with measure, and of great things with weight.

\ciceroSection{102} My whole case for Caecina turned on the words of the interdict: I unraveled tangled matters by defining them, I praised the civil law, I drew distinctions among ambiguous words. In the Manilian law Pompey had to be glorified: with tempered speech I pursued the abundance proper to praise. The whole right of preserving the majesty of the state was bound up in the case of Rabirius: there, accordingly, I blazed forth in every kind of amplification.

\ciceroSection{103} But these things must at times be tempered and varied. What style, then, is not to be found in the seven books of the prosecution? What style is not in the speech for Habitus? In the one for Cornelius? In a great many of my defenses? I would have selected these examples, did I not think they were either well known, or such as those who sought them could read for themselves. For there is no merit of the orator, in any kind, of which there is not in my own speeches at least some specimen — if not the perfection, then still the attempt and the rough sketch.

\ciceroSection{104} We do not attain it; but we see what it is fitting to pursue. For I am speaking now not of myself but of the thing; and in this I am so far from admiring my own work, and am so hard to please and so peevish, that not even Demosthenes himself satisfies me — who, though he alone stands out above all in every kind of speaking, yet does not always fill my ears, so greedy are they and so capacious, and so often do they long for something measureless and boundless.

\ciceroSection{105} But still, since you have come to know this orator thoroughly and most carefully, with his great devotee Pammenes, when you were at Athens, and do not let him slip from your hands — and yet read my works as well — you surely see that he achieves much and we attempt much, that he is able while we are merely willing, to speak in whatever way the case demands. For he, a great man, came after great men himself, and had the greatest orators for his contemporaries; we less so. We should have done a great thing, if indeed we had been able to reach the goal we strive for, in a city in which, as Antonius says, no eloquent man had ever been heard.

\ciceroSection{106} And yet, if Crassus did not seem eloquent to Antonius, or to himself, then Cotta would never have seemed so, nor Sulpicius, nor Hortensius; for Cotta said nothing in the ample style, Sulpicius nothing gently, Hortensius not much with weight; the earlier men were more apt for every kind — I mean Crassus and Antonius. We found, then, that the city's ears were starved of this manifold style of oratory, poured out evenly into every kind, and we were the first — whoever we were, and however little we said — to turn the city to an incredible eagerness for hearing this kind of speech.

\ciceroSection{107} With what loud applause did we, as young men, deliver those words about the punishment of parricides — words which, some while after, we began to feel had by no means cooled enough: ``For what is so common to all as breath to the living, earth to the dead, the sea to those tossed by the waves, the shore to the shipwrecked? They live, while they can, so that they cannot draw breath from heaven; they die so that their bones do not touch the earth; they are tossed by the waves so that they are never washed clean; they are cast up, at the last, so that not even in death do they find rest upon the rocks'' — and what follows; for it is all such as wins a young man praise not so much by its substance and ripeness as by its promise and the expectation it raises. Of the same temper, but already mature, is this: ``the son-in-law's wife, the son's stepmother, the daughter's rival.''

\ciceroSection{108} Nor indeed was this our only ardor, that we should say everything in this manner. For that very youthful exuberance in the speech for Roscius has much that is toned down, and some passages even a little more cheerful — as in the speech for Habitus, in the one for Cornelius, and in several others. For no orator, even in the leisure of Greece, has written so much as my works amount to, and these very works have that variety which I approve.

\ciceroSection{109} Or should I grant to Homer, to Ennius, to the rest of the poets and especially the tragedians, that they need not use the same intensity in every passage but might change it often, sometimes even drawing near to the everyday manner of speech — and yet myself never depart from that fiercest intensity? But why bring forward the poets of divine genius? We have seen actors than whom nothing in their own kind could be more excellent, who not only gave full satisfaction in roles utterly unlike their own, while still keeping within their own line, but whom we have seen win great favor as comic actors in tragedies and tragic actors in comedies: shall I not exert myself?

\ciceroSection{110} When I say ``I,'' I mean you, Brutus; for in my own case what was destined to come about was accomplished long ago. But will you handle every case in the same manner? Or will you reject some kind of case? Or will you maintain, in the same cases, one unbroken and unchanging breath, with no variation at all? Demosthenes, at any rate — whose likeness in bronze I saw lately among the portraits of yourself and your family, when I had come to you at your Tusculan villa (because, I take it, you love him) — yields nothing to Lysias in subtlety, nothing to Hyperides in pointedness and acuteness, nothing to Aeschines in smoothness and splendor of words.

\ciceroSection{111} Many of his speeches are subtle throughout, like the one against Leptines; many wholly weighty, like certain of the Philippics; many varied, like the one against Aeschines on the false embassy, and the one against the same man for the cause of Ctesiphon. As for that middle style, he seizes it whenever he wishes, and, departing from his most weighty manner, it is into this above all that he descends. Yet he stirs up applause, and accomplishes most in speaking, when he uses the passages of weight.

\ciceroSection{112} But let us leave him for a little while, since we are inquiring about the kind, not the man: let us rather unfold the force and nature of the thing itself — that is, of eloquence. Let us nonetheless remember what I said before: that I shall say nothing by way of laying down rules, and shall conduct myself rather so as to seem an appraiser speaking, not a teacher. In this, however, I often go on at greater length, because I see that you are not the only one who will read these pages — you, who hold these matters far better known to you than to us, who seem as it were to be teaching — but that this book, even if less by my recommendation than under your name, must inevitably be made public.

\ciceroSection{113} I hold, then, that it belongs to the perfectly eloquent man not only to possess the faculty proper to him, of speaking copiously and at large, but also to take up that neighboring and bordering science of the dialecticians. And yet oratory seems to be one thing and disputation another, and to converse is not the same as to deliver a speech — and yet both lie in the realm of discourse: let the method of debating and conversing belong to the dialecticians, but that of speaking and adorning to the orators. Zeno, indeed, the founder of the Stoic discipline, used to demonstrate with his hand the difference between these two arts; for when he had closed his fingers and made a fist, he would say that dialectic was of that kind; but when he had drawn them apart and spread out his hand, he would say that eloquence was like that open palm.

\ciceroSection{114} And even before him Aristotle, at the opening of his Art of Rhetoric, says that this art answers, as it were from the other side, to dialectic — differing in this, evidently, that this method of speaking is broader, that one of conversing more compressed. I wish, therefore, that this supreme man should know the whole method of discourse that can be drawn over to the service of speaking; and this discipline — as does not in the least escape you, schooled as you are in these arts — has had a twofold way of teaching. For Aristotle himself handed down a great many precepts of reasoning, and afterward those who are called dialecticians brought forth much that is thornier.

\ciceroSection{115} For my part, I judge that the man who is drawn by the praise of eloquence should not be wholly unschooled in these matters, but trained either in that ancient discipline or in this of Chrysippus. Let him know, first, the force, the nature, and the kinds of words, both simple and combined; then in how many senses each thing is said; by what method it is judged whether a thing be true or false; what follows from each premise, what is consequent upon each thing and what contrary to it; and, since many things are said ambiguously, in what way each of them ought to be divided and explained. These things the orator must master — for they often come up — but, because of themselves they are rather unkempt, a certain polish of speech must be brought to bear in unfolding them.

\ciceroSection{116} And since, in all things that are taught by method and system, one must first establish what each thing is — for unless those who are debating agree on what the matter in dispute is, there can be no right discussion and no coming to an end — our meaning about each thing must often be unfolded in words, and the wrapped-up notion of the thing laid open by definition, if indeed a definition is a statement that shows, as briefly as may be, what that is which is under discussion; then, as you know, once the genus of each thing has been set out, one must see what are the species, or forms, or parts of that genus, so that the whole discourse may be distributed among them.

\ciceroSection{117} There will be in the man we wish to be eloquent, therefore, this faculty: that he can define a thing, yet not do it as tightly and narrowly as is the custom in those most learned discussions, but at once more clearly and also more amply and more accommodated to common judgment and popular understanding; and likewise, when the matter demands, he will partition and divide a whole genus into fixed species, so that none is either passed over or runs to excess. But when he should do this, or in what way, has nothing to do with our present purpose, since, as I said above, I wish to be a judge, not a teacher.

\ciceroSection{118} And indeed let him be furnished not by the dialecticians alone, but let him have all the topics of philosophy known and well handled. For nothing about religion, nothing about death, nothing about piety, nothing about love of country, nothing about good things or evil, nothing about the virtues or the vices, nothing about duty, nothing about pain, nothing about pleasure, nothing about the disorders and errors of the mind — matters that often fall within cases and are pleaded too thinly — nothing, I say, can be spoken and unfolded with weight, with amplitude, with abundance, without that science of which I have spoken.

\ciceroSection{119} I am speaking still of the matter of the speech, not of the kind of speaking itself. For I wish the orator first to have a subject to speak on, one worthy of learned ears, before he considers in what words he is to say each thing, or in what manner; and him too — that he may be the grander, and in a certain way the loftier, as I said above of Pericles — I wish to be not even ignorant of natural philosophy. Surely, when he turns back from celestial matters to human ones, he will both speak and feel more loftily and more magnificently in everything.

\ciceroSection{120} When he has come to know those divine things, I would not have him be ignorant of human things either. Let him have the civil law at his command, of which forensic cases stand in daily need. For what is more disgraceful than to undertake the defense of legal and civil disputes when you are ignorant of the laws and of the civil law? Let him learn, too, the sequence of deeds done and of ancient record — most of all, of course, that of our own state, but also of imperial peoples and illustrious kings; which labor our friend Atticus has lightened for us, who, with the dates preserved and noted, passing over nothing illustrious, has gathered into a single book the record of seven hundred years. But not to know what happened before you were born is to be forever a child. For what is the lifetime of a man, unless it is woven together, through the memory of ancient things, with the age of those who came before? Moreover, the recalling of antiquity and the citing of examples brings to a speech, with the greatest delight, both authority and credit.

\ciceroSection{121} Equipped, then, in this way, he will come to cases, the kinds of which he will first have come to know in themselves. For it will be clear to him that nothing can be in dispute in which the controversy is not raised either by the facts or by the words: by the facts, when it concerns the truth, or the right, or the name of the thing; by the words, when it concerns ambiguity or contradiction. For if at any time one thing seems to lie in the sense and another in the wording, there is a certain kind of ambiguity that commonly arises from a word left unsaid, in which — and this is the very property of ambiguous expressions — we see that two things are signified.

\ciceroSection{122} Since the kinds of cases are so few, the rules for arguments are likewise few. We have been handed two sets of topics from which they may be drawn: the one derived from the facts themselves, the other brought in from outside. It is the handling of the matter, then, that produces an admirable speech; for the facts in themselves lie within an altogether easy understanding. What, after all, now follows that belongs to the art, except: to open the speech in such a way that the hearer is either won over, or roused, or made ready to learn; to set out the matter briefly, plausibly, and clearly, so that what is at issue can be understood; to establish one's own points and overturn the adversary's, and to do this not in confusion but by so concluding the individual arguments that what follows from the premises taken up to prove each point is brought home; and finally, in the peroration, to close by either kindling the feelings or quenching them? How he is to handle each of these parts is hard to say in this place; for they are not always handled in one and the same way.

\ciceroSection{123} But since I am not seeking a man to instruct but a man to approve, I shall give my approval first of all to the one who has seen what is fitting. For this above all is the wisdom the eloquent man must bring to bear: to be the master of occasions and of persons. For I judge that one must not speak always, nor before all audiences, nor against all opponents, nor on behalf of all clients, nor with all associates in one and the same way. He, therefore, will be eloquent who can adapt his speech to whatever is fitting in each case. When he has settled this, then he will speak of each thing as it must be spoken — neither baldly of full matters, nor meagrely of grand ones, nor again the reverse — but his speech will be matched and level with the things themselves.

\ciceroSection{124} His openings modest, not yet inflamed with exalted words, but sharp in their sentiments, whether to wound the adversary or to commend himself. His narratives credible, set out lucidly, not in the language of history but in something close to everyday speech. Then, if the case is slight, the thread of argument too will be slight, both in proving and in refuting, and this will be kept so well in hand that as much accrues to the speech as accrues to the matter.

\ciceroSection{125} But when a case has come his way in which the full power of eloquence can be displayed, then the orator will pour himself out more broadly, then he will rule and bend men's minds and so work upon them as he wills — that is, as the nature of the case and the reason of the occasion shall demand. But that whole admirable adornment, on account of which eloquence has risen to such high honour, will be of two kinds. For while every part of the speech ought to be praiseworthy, so that no word falls from him that is not either weighty or elegant, yet two parts above all are luminous and, as it were, full of action: of these I assign the one to the inquiry into a topic in its universal form, which, as I said above, the Greeks call \textit{[Greek: thesis]}; the other to the heightening and amplifying of things, which by the same men is named \textit{[Greek: auxēsis]}.

\ciceroSection{126} And although this latter ought to be poured out evenly throughout the whole body of the speech, it will nonetheless show to greatest advantage in the commonplaces — called common because they seem to belong alike to many cases, though they ought to be made proper to each single one. But that part of the speech which treats of the universal kind often contains whole cases. For whatever it is in which the contest of the controversy lies, which in Greek is called \textit{[Greek: krinomenon]}, the rule is that it be so stated as to be carried over into the general inquiry, so that one speaks of the universal kind — except when the dispute is about a matter of fact, which is commonly investigated by conjecture.

\ciceroSection{127} It will be argued, however, not after the manner of the Peripatetics — for theirs is an elegant exercise, established long ago, ever since Aristotle — but somewhat more sinewy, and these common matters will be so treated that much is said gently on behalf of the defendants and harshly against the adversaries. As for the heightening of things and, conversely, their disparagement, there is nothing the speech cannot accomplish; and this must be done in the midst of the arguments themselves, as often as occasion is given for amplifying or for diminishing, and almost without limit in the peroration.

\ciceroSection{128} For two things remain which, when well handled by the orator, produce admirable eloquence. The one is what the Greeks call \textit{[Greek: ēthikon]}, suited to natures, to characters, and to the whole habit of life; the other is what those same men name \textit{[Greek: pathētikon]}, by which minds are thrown into turmoil and stirred up, and in this alone the speech reigns supreme. The former is courteous and agreeable, fitted to winning goodwill; the latter is vehement, kindled, driven on, and by it cases are snatched away — and when it sweeps on at speed, it can in no way be withstood.

\ciceroSection{129} In this kind I, though but middling or even much less, yet always wielding great force, have often dislodged my adversaries from their whole footing. On behalf of a friend who stood accused, that supreme orator Hortensius did not answer me; before me, in the Senate, Catiline, that most reckless of men, fell mute under accusation; before me, when in a private case of great gravity Curio the elder had begun to reply, he suddenly sat down, declaring that his memory had been snatched from him by poisons.

\ciceroSection{130} Why should I speak of appeals to pity? Of these I have made the more use because, even when several of us were pleading, the others nevertheless left the peroration to me; and in this I came to seem to excel not by talent but by feeling. Whatever these qualities of mine may be — for I myself am dissatisfied with how great they are — they are at any rate evident in my speeches, though the books lack that breath of life on account of which these same things, when they are pleaded, are wont to seem greater than when they are read.

\ciceroSection{131} And indeed the minds of the jurors must be moved not by pity alone — which I am accustomed to use with such pathos that, while finishing my plea, I once held an infant child in my arms; that in another case, raising up a noble defendant, and lifting up too his little son, I filled the Forum with wailing and lamentation — but it must also be brought about that the juror grows angry and is appeased, that he envies and favours, that he despises and admires, that he hates and loves, that he longs and disdains, that he hopes and fears, that he rejoices and grieves. For this whole range of feelings, the prosecution will furnish examples of the harsher kind, my own defences of the gentler.

\ciceroSection{132} For there is no manner of stirring up or of soothing the mind of a listener that I have not attempted — I would say perfected, did I think so, and did I not, in the face of the truth, dread the charge of arrogance; but, as I said above, it is no force of talent but a great force of feeling that inflames me, so that I cannot contain myself. Nor would the one who listens ever be set alight unless the speech reached him already burning. I would use examples of my own, had you not read them; I would use those of others, whether Latin, if I could find any, or Greek, if it were fitting. But of Crassus there are very few, and those not from the courts; of Antonius nothing, of Cotta nothing, of Sulpicius nothing; Hortensius spoke better than he wrote.

\ciceroSection{133} But how great this power is that we are seeking, let us surmise, since we have no example of it; or, if we are to follow examples, let us take them from Demosthenes — and indeed from a sustained passage of his, from that point where, in the trial of Ctesiphon, he undertook to speak of his own deeds, his counsels, his services to the commonwealth. That speech, surely, can be so brought within that form which is implanted in our minds that no greater eloquence is so much as looked for.

\ciceroSection{134} But now the form itself remains, and that which is called the \textit{[Greek: charactēr]}; and what its character ought to be can be understood from what has been said above. For we have touched on the lights both of single words and of words in combination; in these he will so abound that no word leaves his lips that is not either elegant or weighty, and from every kind there will be very frequent metaphors, because these, by virtue of likeness, transfer the mind and carry it back and move it this way and that — and this motion of thought, swiftly stirred, is of itself a delight. And the remaining devices that are drawn from the arrangement of words, the lights as it were, lend great adornment to a speech; for they are like those things which, in the rich furnishing of a stage or of the Forum, are called showpieces — not because they alone adorn, but because they stand out.

\ciceroSection{135} The same holds of those things which are the lights of a speech and in a manner its showpieces: when words are doubled and repeated, or set down slightly altered; or when the speech is led off more than once from the same word, or is gathered back into the same word, or both; or when the same word is joined on in repetition, or the same word is carried to the end; or when one and the same word is set down continuously but not in the same sense; or when words fall alike or end alike; or when contraries are matched against contraries; or when one mounts step by step upward; or when, with the conjunctions stripped away, several things are said in loose succession; or when, passing something over, we show why we do so; or when we correct ourselves as if reproaching ourselves; or if there is some exclamation, whether of wonder or of complaint; or when the cases of the same noun are shifted again and again.

\ciceroSection{136} But the ornaments of the thoughts are greater; and because Demosthenes uses these most frequently, there are those who think his eloquence is for that reason supremely praiseworthy. And indeed scarcely any passage is spoken by him without some shaping of the thought; nor is anything else meant by speaking than to illuminate all, or at any rate most, of the thoughts with some figure. Since you, Brutus, have an excellent grasp of these, what need is there to use the names or the examples? Let the topic merely be marked.

\ciceroSection{137} So, then, the man we are seeking will speak in such a way as to turn the same thing over often in many ways, and to dwell on one point and linger over the same thought; often too to make light of something, often to mock it; to swerve from his theme and bend the thought aside; to announce what he is about to say; to define, when he has now finished some matter; to call himself back; to repeat what he has said; to round off an argument with reasoning; to press by questioning; and again, as if in answer to the questions, to answer himself; to wish a thing taken and felt the contrary of what he says; to raise a doubt whether there is anything he should rather say, or how he should say it; to divide into parts; to leave something aside and pass it by; to fortify himself beforehand; on the very point for which he is reproached, to lay the fault on the adversary;

\ciceroSection{138} often to deliberate, as it were, with those who are listening, and sometimes even with the adversary; to describe the speech and the characters of men; to bring in mute things speaking; to turn the mind away from the matter at hand; often to turn it to mirth and laughter; to seize beforehand on what seems likely to be raised against him; to draw comparisons; to use examples; to apportion, assigning one thing to one and another to another; to check an interrupter; to say that he passes something over in silence; to give warning of what they should beware; to dare something rather freely; even to grow angry, and at times to rebuke; to entreat, to beseech, to provide a remedy; to swerve somewhat from his theme; to wish, to curse; to make himself intimate with those before whom he will speak.

\ciceroSection{139} And he will pursue still other excellences, as it were, of speaking: brevity, if the matter calls for it; often, too, by his words he will set the thing before the eyes; often he will carry it beyond what could really happen; the suggestion will often be greater than the words; often there will be mirth, often the mimicry of life and of characters. In this kind — for you see something like a forest — the whole greatness of eloquence ought to shine forth.

\ciceroSection{140} But these things, unless they are arranged and, as it were, built up and bound together by the words, cannot attain to the distinction we wish for. And although, when I saw that I must speak of this next, I was already moved by those considerations I set out above, I was nonetheless more troubled by what follows. For it occurred to me that one might find not only the envious, of whom all places are full, but even the supporters of my own renown, who would not think it the part of a man — about whose services the Senate had passed such judgments, with the Roman people's approval, as it had passed about no one — to commit so many things to writing about the craft of speaking. To these, if I gave no other answer than that I had been unwilling to refuse Marcus Brutus when he asked, the excuse would be just, since I wished to satisfy a man both most dear to me and most distinguished, asking for what was right and honourable.

\ciceroSection{141} But if I were to declare — and would that I could! — that I would hand on to those eager for speaking the precepts and, as it were, the roads that lead to eloquence, who, after all, that is a just judge of things would find fault with it? For who has ever doubted that, in our commonwealth, eloquence has always held the first place in the affairs of the city and of peace, and the knowledge of the law the second? Since in the one there was the greatest store of favour, of glory, of protection, while in the other there was the laying down of formulas and securities — which itself often sought aid from eloquence, and against its opposition could scarcely defend its own borders and bounds.

\ciceroSection{142} Why, then, has it always been honourable to teach the civil law, and have the houses of the most illustrious men flourished with pupils — yet, if anyone sharpens or aids the young in speaking, is he to be reproached? For if it is a fault to speak with ornament, let eloquence be driven out of the state altogether; but if it not only adorns those who possess it but even aids the whole commonwealth, why is it base to learn what it is honourable to know, or why is it not glorious to teach what it is most splendid to be able to do?

\ciceroSection{143} `But the one has been long practised, the other is new.' I grant it; but there is a reason for both. For it was enough to hear the one set of men giving their responses, so that those who taught set aside no time of their own for the purpose, but at one and the same time satisfied both their students and their consultants; whereas the other set, spending their household hours in studying and composing cases, their hours in the Forum in pleading them, and what remained in restoring themselves — what room did they have for instructing or teaching? And I rather think that most of our orators, unlike myself, have been stronger in talent than in learning; and so they could perhaps speak better than they could teach, and I the reverse.

\ciceroSection{144} `But teaching carries no dignity.' Certainly not, if it is done as in a school; but if by advising, by exhorting, by inquiring, by sharing, and at times even by reading and listening together — I do not see why, if by teaching as well you can sometimes make men better, you should be unwilling to do so. Or is it honourable to teach the words by which a transfer of sacred rites is effected — as indeed it is — and not honourable to teach the words by which the sacred rites themselves may be retained and defended?

\ciceroSection{145} `But men profess the law even when they do not know it; whereas eloquence even those who have attained it nevertheless pretend they have no power in.' This is because good sense is welcome to men, but the tongue is suspect. Can eloquence, then, lie hidden? Or does what it dissembles escape notice? Or is there any danger that someone should think it base, in a great and glorious art, to teach others that which it was most honourable for oneself to have learned?

\ciceroSection{146} And perhaps the rest are more guarded; I have always made open profession that I learned. For how could I — I, who in my youth had gone away from home, and had crossed the seas for the sake of these studies, and whose house was full of the most learned men, and in whose conversation there were perhaps some marks of learning, and whose writings were widely read by the public — how could I pretend that I had not learned? What reason was there for me to win approval, except that I had perhaps made too little progress? And though this is so, yet those matters of which I spoke above had more dignity, in the discussion, than these of which I am now to speak.

\ciceroSection{147} For we shall speak of the composing of words and of the all but counting and measuring out of syllables; and these, even if they are necessary, as they seem to me to be, are nonetheless practised more magnificently than they are taught. This is altogether true, but it is said with particular aptness in the present matter. For in all great arts, as in trees, it is the height that delights us, not in the same way the roots and the stocks; yet the former cannot exist without the latter. As for me, whether that most familiar verse, which forbids us to be ashamed to profess openly the art we practise, does not allow me to dissemble the thing in which I take delight; or whether your own zeal has wrung this volume from me — still, an answer had to be given to those whom I suspected would find something to reproach.

\ciceroSection{148} But even if what I have said were not so, who, after all, would show himself so hard and boorish as not to grant me this indulgence: that, when my arts of the Forum and my public pleadings had collapsed, I should give myself not to idleness, which I cannot do, nor to gloom, which I resist, but rather to letters? — letters which formerly led me into the courts and into the Senate-house, and now console me at home; and not only with such matters as this book contains, but with far weightier and greater ones as well. If these are brought to completion, then surely my domestic, indoor letters will answer to my greatest affairs of the Forum, both at home and abroad. But let us return to the discussion we have undertaken.

\ciceroSection{149} Words, then, will be arranged either so that they cohere among themselves as aptly as possible, last with first, and so that they consist of the most pleasant-sounding utterances; or so that the very form and symmetry of the words completes its own rounded period; or so that the whole construction falls rhythmically and aptly. And let us look first at the nature of that which most of all demands diligence: that there be formed, as it were, a certain structure, and yet that it not be formed with visible labour; for the toil would be both endless and childish. This is what Scaevola, in Lucilius, wittily mocks in Albucius: `How prettily your phrases are set together, like little tesserae, all in mosaic-work, with inlaid patterns running like worm-tracks!'

\ciceroSection{150} I do not wish this minute construction to show itself; yet a practised pen will readily produce a formula for composing. For just as in reading the eye, so in speaking the mind will look ahead to what follows, lest the collision of the last words with the first that come after produce either gaping or harsh sounds. For however pleasant and weighty the thoughts may be, nonetheless, if they are uttered with the words ill-set, they will offend the ears, whose judgment is most haughty. And this, indeed, the Latin tongue so heeds that no one is so rustic as to be unwilling to join the vowels.

\ciceroSection{151} On this point some even find fault with Theopompus for having so studiously avoided that collision of vowels, although his own teacher Isocrates had done the same. But Thucydides did not avoid it, nor did Plato, a writer not a little greater than he — and not only in those compositions which are called dialogues \textit{[Greek: dialogoi]}, where it had even to be done on purpose, but also in that public oration in which it is the custom at Athens to praise before the assembly those who have fallen in battle: an oration so highly approved that, as you know, it must be recited every year on the appointed day. In it that frequent meeting of vowels occurs, the very thing which Demosthenes for the most part avoids as a fault.

\ciceroSection{152} But the Greeks may see to their own affairs; for us it is not permitted, even if we should wish it, to pull voices apart. Those somewhat rough orations of Cato bear witness to it; all the poets bear witness, except those who, to make their verse come out, often left a gap of vowels, as Naevius does — ``you who dwell along the Danube's stream and the cold,'' and likewise: ``which never to you Greeks and barbarians''; whereas Ennius does so often — ``unconquered Scipio'' — and once at least — ``with this motion the radiant Etesian winds upon the shoals of the sea.''

\ciceroSection{153} This same thing our writers would more often not have tolerated, which the Greeks are even accustomed to praise. But why do I speak of vowels? Even without vowels they often contracted words for the sake of brevity, so that they would say \textit{multi' modis} for ``in many ways,'' \textit{in vas' argenteis} for ``in vessels of silver,'' \textit{palmi' crinibus} for ``with palm-like locks,'' \textit{tecti' fractis} for ``with shattered roofs.'' And what could be more licentious than that they even contracted the names of men, so that these might be the more fitting? For just as \textit{duellum} became \textit{bellum} and \textit{duis} became \textit{bis}, so the man who crushed the Carthaginians at sea — Duellius — they named Bellius, although his ancestors had always been called Duellii. Indeed words too are often contracted not for use's sake but for the ear's. For how did your own Axilla become Ala, if not by the flight of a too-rugged letter? — a letter which the elegant usage of the Latin tongue plucks out also from \textit{maxillae} and \textit{taxilli} and \textit{paxillus} and \textit{vexillum}.

\ciceroSection{154} They also gladly joined words by coupling them, as \textit{sodes} for \textit{si audes} (``if you please''), \textit{sis} for \textit{si vis} (``if you wish''). And in the single word \textit{capsis} there are three words. They say \textit{ain} for \textit{aisne}, \textit{nequire} for \textit{non quire}, \textit{malle} for \textit{magis velle}, \textit{nolle} for \textit{non velle}; and we often say \textit{dein} and \textit{exin} for \textit{deinde} and for \textit{exinde}. And tell me, does it not betray its own origin, the fact that we say \textit{cum illis} but do not say \textit{cum nobis}, but rather \textit{nobiscum}? — because, if it were said the other way, the letters would meet too obscenely, as in fact they would have met just now, had I not interposed a word. From this comes \textit{mecum} and \textit{tecum}, not \textit{cum me} and \textit{cum te}, so that it should be like \textit{nobiscum} and \textit{vobiscum}.

\ciceroSection{155} And there are even some who now, all too late, set about correcting antiquity, finding fault with these usages. For in place of \textit{deum atque hominum fidem} (``by the faith of gods and men'') they say \textit{deorum}. So, I suppose, the old writers did not know this? Or did custom grant them this license? And so that same poet who had made these more unusual contractions — ``of my father I am ashamed of my deed'' (\textit{patris mei meum factum pudet}) for \textit{meorum factorum}, and ``it is woven, ruin sweeps off the swarm'' for \textit{exitiorum} — does not say \textit{liberum}, as most of us speak when we say ``fond of children'' or use the word in the genitive plural of children, but as those archaizers prefer: ``nor may you ever lift into your lap a brood of children sprung from you'' (\textit{liberorum}), and again: ``for of Aesculapius' children'' (\textit{liberorum}). But that other poet, in his \textit{Chryses}, gives us not only — ``citizens, ancient friends of my forefathers'' (\textit{maiorum meum}), which was customary — but even, more harshly: ``partners in counsel, interpreters of augury and entrails'' (\textit{socii}, \textit{augurium atque extum}); and the same goes on: ``after the horror-bearing prodigy, the portent of dread'' (\textit{pavos}) — forms which are surely not in common use among all neuter nouns. For I would not so gladly say \textit{armum iudicium} (``the judgment of arms'') — though it stands in that same poet: ``has nothing reached you about the judgment of the arms?'' — as I would say ``a century of carpenters'' (\textit{centuriam fabrum}) and ``of suitors'' (\textit{procum}),

\ciceroSection{156} as the censors' registers express it — I make bold to say it — not \textit{fabrorum} or \textit{procorum}; while plainly ``the court of the two men'' (\textit{duorum virorum}) or ``of the three men for capital cases'' (\textit{trium virorum}) or ``of the ten men for adjudicating disputes'' (\textit{decem virorum}) I say never. And what did Accius say? ``I see two tombs of two bodies'' (\textit{duorum corporum}); and the same poet, ``one woman of two men'' (\textit{duum virum}). I understand what is correct; but sometimes I speak as is conceded, as I say either \textit{pro deum} or \textit{pro deorum}, and at other times as is necessary, when I say \textit{trium virum}, not \textit{virorum}, and \textit{sestertium nummum}, not \textit{sestertiorum nummorum}, since in these cases usage is not divided.

\ciceroSection{157} What of the fact that they forbid us to say \textit{nosse}, \textit{iudicasse}, and bid us say \textit{novisse} and \textit{iudicavisse}? As if, in fact, we did not know that in this kind both the full word is rightly said and the shortened one is said by usage! And so Terence uses both: ``Hey, you — did you not know your own kinsman?'' (\textit{non noras}); and afterward the same poet, ``Stilpo, I say — you had known him'' (\textit{noveras}). \textit{Sient} is the full form, \textit{sint} the shortened one; one may use either. So in the same author: ``how dear these things are they perceive only by lacking them, and how great the dominion of the things they hold'' (\textit{sient}). Nor indeed would I reproach \textit{scripsere alii rem}, although I feel that \textit{scripserunt} is the truer form; but I gladly defer to a usage that indulges the ears. ``The same plain holds it,'' says Ennius (\textit{isdem}); and in the temples, EIDEM PROBAVIT (``the same men approved it''); yet \textit{isdem} was truer, and not, however, \textit{eisdem}, which would be too rich; \textit{isdem} sounded ill; from custom it was obtained that one might err for the sake of euphony. And I would more gladly say \textit{posmeridianas quadrigas} than \textit{postmeridianas quadriiugas}, and, by Hercules, \textit{me hercule} rather than \textit{me hercules}. \textit{Non scire} now seems quite barbarous; \textit{nescire} is sweeter. And the very word \textit{meridies} (``midday'') — why not \textit{medidies}?

\ciceroSection{158} Because, I suppose, it was less pleasant. There is the one preposition \textit{af}, which now survives only in the ledgers of receipts — and not even in those of everyone — while in the rest of speech it has been altered; for we say \textit{amovit} and \textit{abegit} and \textit{abstulit}, so that by now you can hardly tell whether the true form is \textit{a}, or \textit{ab}, or \textit{abs}. And what of this — that \textit{abfugit} seemed ugly and they would not have \textit{abfer}, but preferred \textit{aufugit} and \textit{aufer}? — a preposition which, except in these two words, will be found in no other word at all. \textit{Noti} and \textit{navi} and \textit{nari} were known forms; and since \textit{in} ought to be prefixed to them, it seemed sweeter to say \textit{ignotos}, \textit{ignavos}, \textit{ignaros} than as truth demanded. They say \textit{ex usu} and \textit{e re publica}, because in the one a vowel followed, while in the other there would be harshness unless you removed the letter, as in \textit{exegit}, \textit{edixit}; \textit{refecit}, \textit{rettulit}, \textit{reddidit}: the first letter of the joined word altered the preposition, as in \textit{subegit}, \textit{summovit}, \textit{sustulit}.

\ciceroSection{159} And what of joined words? How neatly we say \textit{insipientem}, not \textit{insapientem}; \textit{iniquum}, not \textit{inaequum}; \textit{tricipitem}, not \textit{tricapitem}; \textit{concisum}, not \textit{concaesum}! On the strength of which some would even have \textit{pertisum}, which that same custom did not approve. And what could be more refined than this, which arises not by nature but by a kind of established practice? We say \textit{indoctus} with the first letter short, \textit{insanus} with it lengthened; \textit{inhumanus} short, \textit{infelix} long. And, not to be tedious, in those words whose opening letters are the same as in \textit{sapiens} and \textit{felix}, it is pronounced long; in all the rest, short; and likewise \textit{composuit}, \textit{consuevit}, \textit{concrepuit}, \textit{confecit}. Consult truth: it will rebuke you; refer the matter to the ears: they will approve. Ask why it is so: they will say, because it gives pleasure. And speech ought to gratify the pleasure of the ears.

\ciceroSection{160} Indeed I myself, although I knew that our ancestors had so spoken as to use the aspirate nowhere except before a vowel, used to speak in such a way as to say \textit{pulcros}, \textit{Cetegos}, \textit{triumpos}, \textit{Cartaginem}; at length, and that late, when the truth had been wrung out of me by the ears' protest, I yielded the usage of speech to the people and kept the knowledge to myself. Yet we say \textit{Orcivii} and \textit{Matones}, \textit{Otones}, \textit{Caepiones}, \textit{sepulcra}, \textit{coronas}, \textit{lacrimas}, because the judgment of the ears permits it. Ennius always says Burrus, never Pyrrhus; ``by force they laid open'' — Bruges, not Phryges — as his own ancient books declare. For they did not employ the Greek letter, whereas now they employ even two; and since one would have to say \textit{Phrygum} and \textit{Phrygibus}, it would be absurd either to employ a Greek letter even in barbarian case-endings or to speak Greek only in the nominative; yet we say both Phryges and Pyrrhus for the ears' sake.

\ciceroSection{161} Indeed there is something else — by now it seems rather countrified, though in old times it was more polished. Of those words whose last two letters were the same as in \textit{optimus}, they used to drop the final letter, unless a vowel followed. So there was not in their verses that offense which the modern poets now avoid. For we used to speak thus: ``who is foremost of all'' (\textit{omnibu' princeps}), not \textit{omnibus princeps}; and ``worthy of that life and place'' (\textit{dignu' locoque}), not \textit{dignus}. But if untaught custom is so skilled an artificer of euphony, what should we suppose may be demanded of art itself and of learning?

\ciceroSection{162} I have said this more briefly than if I were arguing on this single subject — for this is a topic of wide extent, concerning the nature and the use of words — yet at greater length than my settled plan required. But since the judgment of matters and of words rests in good sense, while of sounds and rhythms the ears are the judges; and since the former are referred to the understanding, the latter to pleasure; and since in the former reason discovers the art, in the latter the senses do — it followed that either we had to neglect the inclination of the ears, by which we strove to win approval, or else the art of conciliating them had to be found out.

\ciceroSection{163} There are, then, two things that soothe the ears: sound and rhythm. Of rhythm presently; now we inquire about sound. Words, as we said above, are to be chosen above all for sounding well — yet not, as the poets do, sought out for their sound, but taken from the common stock. ``Where the sea of Helle, beyond Tmolus and the Tauric peaks'' — the verse is illuminated by the splendid names of places, but the next is befouled by a most unpleasant letter: ``the boundary of fruitful Asia, and her crowded fields, it holds.''

\ciceroSection{164} Therefore let us use the soundness of our own words rather than the splendor of Greek ones — unless perhaps we are loath to speak thus: ``at the season when Paris carried off Helen,'' and what follows. No, rather let us follow words like these, and shun harshness: ``I hold that creature of dreadful crashing din,'' and likewise: ``crooked-spoken malice.'' Nor will words be arranged by reason only, but they will also be brought to a close — since we have said that this is the second judgment of the ears. And they are brought to a close either by the arrangement itself and, as it were, of their own accord, or by a certain kind of words which themselves contain a symmetry. These — whether they have like case-endings at the close, or like is matched with like, or contraries are set against each other — are by their own nature rhythmical, even if nothing has been done on purpose.

\ciceroSection{165} In the pursuit of this symmetry we have been told that Gorgias was the leader; of which kind are those words of ours in the speech for Milo: ``For this, members of the jury, is a law not written but inborn — one which we have not learned, received, or read, but which we have seized, drawn off, and pressed out from nature herself; for which we were not trained but made, not instructed but steeped.'' For things of this sort are such that, because they are referred to the end to which they ought to be referred, we recognize that the rhythm was not sought after but followed of itself.

\ciceroSection{166} The same happens in the matching of contraries, as in those phrases by which not only rhythmical prose but even verse is produced: ``her whom you accuse of no wrong, you condemn'' — he would have said \textit{condemn} (\textit{condemnas}) who wished to escape a verse — ``and her whom you declare to have deserved well, do you say deserves ill? what you know profits nothing; what you do not know does harm?'' The very matching of contraries produces a verse. The same would be rhythmical in prose: ``what you know profits nothing; what you do not know does much harm.'' Always these figures, which the Greeks call antitheses \textit{[Greek: antitheta]}, when contraries are set against contraries, produce an oratorical rhythm by the very necessity of the thing, even without effort.

\ciceroSection{167} The ancients took delight in this kind even before Isocrates, and most of all Gorgias, in whose discourse the symmetry itself for the most part produces the rhythm. We too have used it frequently in this kind, as in those passages in the fourth book of the prosecution: ``Compare this peace with that war, the coming of this praetor with the victory of that commander, the foul retinue of this man with the unconquered army of that one, the lusts of this man with the self-command of that one: from him who took it you will say Syracuse was founded; from this man who received it founded, that it was taken.''

\ciceroSection{168} So let these rhythms be acknowledged, and let that third kind be unfolded — what its nature is, the nature of rhythmical and well-fitted prose. Those who do not perceive it, what sort of ears they have, or what there is in them resembling a human being, I do not know. My own, at any rate, take joy in a perfected and completed circuit of words, and feel a shortfall, and have no love for things that run over. And why do I say my own? I have often seen assemblies cry out aloud when the words have fallen aptly together. For this is what the ears await — that the thought be bound up by the words. ``This was not so among the ancients.'' And indeed it was almost the only thing that was not so among them; for they both chose their words and devised thoughts weighty and pleasing, but they bound them, or filled them out, too little.

\ciceroSection{169} ``This is the very thing that delights me,'' they say. But what if that most ancient kind of painting, with its few colors, should please us more than this kind, now brought to perfection — must we, I suppose, go back to the one and reject the other! They glory in the names of the ancients. But just as in human ages old age has authority, so in models antiquity has it — and this carries the greatest weight with me too. Nor do I demand what antiquity lacks rather than praise what it has, especially since I judge the things it has to be greater than those it lacks. For there is more of value in words and in thoughts, in which those men excel, than in the rounding-off of thoughts, which they do not have. The rounded close was discovered later — which, I believe, those men of old would have used, had the thing then been known and in practice; and once it was discovered, we see that all the great orators used it.

\ciceroSection{170} But the name carries odium, when in judicial and forensic oratory rhythm is said to be present — \textit{numerus} in Latin, \textit{[Greek: rhuthmos]} in Greek. For it seems that too much of a snare is laid for catching the ears, if rhythms too are sought by the orator in his speaking. Trusting in this, those men themselves speak in a broken and clipped fashion, and abuse those who deliver in a well-fitted and rounded way: if the words are empty and the thoughts trivial, with justice; but if the matter is sound and the words well chosen, what reason is there why they should prefer to make their speech limp or come to a halt rather than run on level with the thought? For this much-resented rhythm brings nothing else than this — that the thought be aptly compassed by the words; which is done even by the ancients, but for the most part by chance, often by nature; and the passages of theirs that are greatly praised are praised, as a rule, just because they are rounded off.

\ciceroSection{171} And among the Greeks indeed it is now nearly four hundred years since this has been approved; we have only lately come to acknowledge it. Therefore, if it was permitted to Ennius, in his contempt for the old, to say, ``in the verses which once the Fauns and seers used to chant,'' shall it not be permitted me to speak in the same way of the ancients? — especially since I am not going to say, as he does, ``before this man,'' nor what follows: ``I have dared to unbar.'' For I have both delivered speeches and heard not a few in which the prose was brought to an all but perfect close. Those who cannot do this are not content not to be despised; they wish even to be praised. As for me, I praise those very men whose imitators these people call themselves — and praise them deservedly, even though I miss something in them; but these others not at all, who follow nothing of those men except their faults, while they fall as far as possible short of their merits.

\ciceroSection{172} But if they have ears so uncultivated and so boorish, will not even the authority of the most learned men move them? I pass over Isocrates and his disciples Ephorus and Naucrates — although, as the most weighty authorities on composing and adorning a speech, they ought themselves to have been the foremost orators. But who of all men was more learned, who keener, who more incisive in either inventing or judging matters, than Aristotle? And who, besides, opposed Isocrates more fiercely? Yet that very man forbids verse to be in oratory and bids there be rhythm. His pupil Theodectes — a writer and craftsman of the first polish, as Aristotle often signifies — both holds and teaches this same thing; and Theophrastus indeed treats the same subjects even more carefully. Who, then, would endure these people, who do not approve of such authorities? — unless they are altogether ignorant that these precepts come from them.

\ciceroSection{173} But if that is so — and I judge it to be no otherwise — well then, are they not moved by their own senses? Does nothing seem empty to them, nothing unshaped, nothing curtailed, nothing limping, nothing overflowing? In verse, at any rate, whole theaters cry out if there has been a single syllable too short or too long; and yet the crowd knows nothing of metrical feet, holds no rhythms, and does not understand what offends it, or care, or grasp wherein the offense lies; and nonetheless nature herself has lodged in our ears the judgment of all lengths and shortnesses in sounds, just as of high and low pitches of voice.

\ciceroSection{174} Do you wish then, Brutus, that we should unfold this whole topic even more carefully than those very men did who handed down both this and other matters to us — or can we be content with what has been said by them? But why do I ask whether you wish it, when from your most learnedly written letter I have perceived that you wish it above all else? First, then, let the origin be unfolded; next the cause; then the nature; and last of all the practice itself of well-fitted and rhythmical prose. For those who most admire Isocrates count it among his highest praises that he was the first to join rhythms to unbound prose. For when he saw that orators were heard with severity but poets with pleasure, then, it is said, he sought after rhythms, to use them even in prose, both for the sake of charm and so that variety might forestall satiety.

\ciceroSection{175} Which is said by them with some truth in part, but not as a whole. For it must be granted that no one handled this kind more knowingly than Isocrates; but the first to discover it was Thrasymachus, all of whose writings survive — and indeed survive too rhythmically. For, as I said a little before, the joining of like to like and the closing of clauses in like manner, and likewise the matching of contraries with contraries — which of their own accord, even if you do not work at it, for the most part fall out rhythmically — Gorgias was the first to discover, but he used them too immoderately. And this kind, as was said before, is one of the three parts of arrangement.

\ciceroSection{176} Both of these men came before Isocrates in age, so that he surpassed them in restraint, not in invention. For as he is more temperate in transferring and coining words, so is he more sedate in the rhythms themselves. But Gorgias is too greedy for that kind, and abuses these elegant flourishes — for so he himself esteems them — with too little restraint; which Isocrates nonetheless tempered more moderately, after he had heard Gorgias as a young man in Thessaly, when the man was already old. Indeed, the further he advanced in age — for he completed nearly a hundred years — the more he relaxed his hold from the excessive constraint of rhythms, as he declares in that book which he wrote to Philip of Macedon, when he was by then quite an old man; in which he says that he now serves rhythms less than he had been wont. So he had corrected not only his predecessors but even his own self.

\ciceroSection{177} Since, then, we have those leaders and originators of well-fitted prose whom we have named, and the origin has been found, let the cause be sought. It is so plain that I wonder the ancients were not moved by it — especially since, as happens, they often by chance said something in a closely rounded and well-fitted way. And when this had struck the minds and the ears of men, so that one could perceive that what chance had poured out had fallen out pleasantly, surely the kind was to be marked, and they ought to have imitated themselves. For the ears themselves, or the mind through the report of the ears, contain within themselves a certain natural measure of all utterances.

\ciceroSection{178} And so they judge both the longer and the shorter, and always look for the perfect and the measured; some things they feel to be mutilated and, as it were, cut short, and at these they take offense, as though defrauded of what is owed; others they feel to be too drawn out and, as it were, running over with too little restraint, and these the ears reject even more — as in most things, so in this kind, what is excessive offends more violently than what seems too little. Therefore, just as the verse of poetry was discovered by the ears' setting of bounds and by the observation of skilled men, so in prose it was noticed — much later, to be sure, but with the same nature giving the prompt — that there are certain fixed courses and closes of words.

\ciceroSection{179} Since, then, we have shown the cause too, let us now, if you please, unfold the nature — for that was the third point; a discussion which does not belong to the plan of this present conversation, but to the innermost art. For it may be asked what the rhythm of prose is, and where it is placed, and from what it has arisen, and whether it is one or two or more, and by what method it is composed, and to what end and when and in what place and in what manner it brings, when employed, some measure of pleasure.

\ciceroSection{180} But as in most matters, so in this, there is a twofold way of considering it: of which the one is longer, the other shorter — and the same also plainer. The first question belonging to the longer way is whether there is any rhythmical prose at all; for to some it does not seem so, because there is nothing fixed in it as there is in verse, and because those very men who affirm that these rhythms exist cannot render the reason why they exist. Next, if there is rhythm in prose, of what sort it is or sorts, and whether from the poetic rhythms or from some other kind, and, if from the poetic, which of them it is or are; for to some only one, to others several, to others all of them seem the same. Next, whatever they are, whether one or several, whether they are common to every kind of prose — since one kind is for narrating, another for persuading, another for instructing — or whether different rhythms are suited to each kind of prose; if common, which they are; if different, what the difference is, and why rhythm does not appear in prose as equally as in verse.

\ciceroSection{181} Next, when we call a piece of prose rhythmical, is that quality produced by rhythm alone, or also by a certain arrangement, or by a certain choice of words? Or has each its own province — rhythm appearing in the intervals, arrangement in the sounds, the very choice of words showing itself as a kind of form and brilliance in the discourse — so that arrangement is the source of all, and from it are produced both the rhythm and those things that are called, as it were, the forms and brilliances of speech, which, as I said, the Greeks call \textit{[Greek: schēmata]} (figures)?

\ciceroSection{182} But it is not one and the same thing that is pleasing to the ear, that is brought to perfection by measure, and that is illuminated by the choice of words — though this last is indeed close kin to rhythm, since for the most part it is complete in itself. Arrangement, however, differs from both, for it is wholly given over to the weight or the sweetness of the sounds. These, then, are roughly the matters in which the nature of the thing must be sought.

\ciceroSection{183} That there is, then, a certain rhythm in prose is not hard to recognise. For the ear is the judge of it; and in this it is unfair not to acknowledge what occurs simply because we cannot discover why it occurs. For the line of verse itself was not recognised by reasoning, but by nature and by the ear, which reasoning, having measured it out, afterward taught what was happening. Thus the observation of nature and close attention gave birth to the art. But in verse the matter is more open to view — although prose too, when certain measures are stripped of their song, seems to be set free into rhythm, and most of all in each of the best of those poets whom the Greeks name \textit{[Greek: lyrikoi]} (lyric poets): once you have stripped these of their song, what remains is almost bare prose.

\ciceroSection{184} There are some things much like these even among our own poets, such as that passage in the \textit{Thyestes}: ``Whom shall I say you are? you who in halting old age \dots'' — and what follows; which, unless the flute-player has joined in, are most like prose set loose. But the senarii of the comic poets, because of their resemblance to ordinary talk, are often so flattened that sometimes the rhythm and the verse can scarcely be made out in them. For which reason rhythm is harder to find in prose than in verse.

\ciceroSection{185} In all there are two things that season prose: the pleasantness of the words and of the rhythms. In the words there lies, as it were, a kind of raw material, but in the rhythm the polishing. Yet, as in everything else, the inventions of necessity are older than those of pleasure.

\ciceroSection{186} And so both Herodotus and the same and the earlier age were without rhythm, except where it came in at random and by chance; and the very ancient writers left us nothing at all about rhythm, but many precepts about prose — for what is both easier and more necessary is always learned first. Thus words transferred in meaning, or coined, or compounded were easily understood, because they were taken from common and everyday speech. Rhythm, on the other hand, was not fetched from the household stores, nor did it have any necessary tie or kinship with prose. And so, observed and recognised somewhat later, it brought to prose a kind of training-ground and its final outlines.

\ciceroSection{187} But if one kind of prose is cramped and clipped, and another expanded and flowing, this must come about not from the nature of the letters but from the variety of long and short intervals; and since prose, interwoven and blended of these, is now firm-footed, now swift in motion, it follows that the nature of rhythm is contained in a force of this kind. For that periodic sweep, which we have now often spoken of, is carried along and glides the more rapidly for the rhythm itself, until it reaches its end and comes to rest. It is clear, then, that prose ought to be bound by rhythms and free of verses.

\ciceroSection{188} But whether these rhythms are poetic or of some other kind is the next thing to be examined. There is, then, no rhythm outside the poetic, for the kinds of rhythm are fixed and limited. For every such rhythm is of a sort to be one of three. The foot, which is brought into use for rhythm, is divided into three parts, so that one part of the foot must be either equal to the other part, or twice as great, or half as great again. Thus the dactyl turns out equal, the iamb double, the paean half again as great; and how can these feet fail to fall into prose? When they are set in due order, what results must of necessity be rhythmical.

\ciceroSection{189} But the question is what rhythm, or which rhythms above the rest, are to be used. That they all fall into prose can be understood from this too — that we often utter verses in prose without meaning to. This is a serious fault, but we do not attend to it nor catch ourselves at it; senarii indeed, and Hipponactean lines, we can scarcely escape, for our prose consists in great part of iambs. And yet the hearer easily recognises those verses, for they are the most familiar; but unawares we often work in even less familiar ones, which are verses all the same — a faulty thing, and to be guarded against by long forethought of the mind.

\ciceroSection{190} Hieronymus the Peripatetic, a man of the first distinction, picked out of the many books of Isocrates perhaps thirty verses, most of them senarii, but anapaests too — and what could be more disgraceful than that? Yet in the selecting he acted with malice; for, taking away the first syllable from the first word of a sentence, he joined to the last word again the first syllable of the following sentence; and so was produced the anapaest that is called Aristophanic. That this should not happen can neither be watched against nor needs to be. But for all that, this fault-finder, in the very passage in which he finds fault — as was noticed by me when I was studiously inquiring into his case — unawares lets slip a senarius of his own. Let this, then, be taken as established: that rhythms are present even in unbound prose, and that they are the same for the orator as for the poet.

\ciceroSection{191} It follows, then, that we must consider which rhythms fall best into a well-shaped prose. For there are those who think it is the iambic, since it is most like ordinary prose — and that this is why it is the foot chiefly employed in plays, because of its likeness to real speech, whereas that dactylic rhythm of the hexameter is better suited to grand utterance. Ephorus, however — himself a smooth orator, and trained in the best school — pursues the paean or the dactyl, but shuns the spondee and the trochee. Because the paean had three short syllables, and the dactyl two, he holds that the brevity and quickness of the syllables makes the words glide more readily, and that the opposite happens with the spondee and the trochee: since the one consists of long syllables and the other of short, the prose becomes either too headlong or too slow, and neither is held in measure.

\ciceroSection{192} But both those earlier men are mistaken, and Ephorus is at fault. For those who pass over the paean do not see that they are passing over the gentlest of rhythms, and at the same time the most ample. Aristotle holds quite otherwise: he judges that the heroic rhythm is grander than unbound prose requires, while the iamb is drawn too much from common talk. So he approves a prose neither low and abject nor too lofty and overblown, but he would have it full of weight, that it may be able to carry its hearers over into greater admiration.

\ciceroSection{193} The trochee, however, which occupies the same span as the choree, he calls the \textit{cordax} (the lewd dance-step), because its contraction and brevity have no dignity. Thus he approves the paean, and says that all use it, but that they do not feel it when they do; that it is a third kind, midway between those others; and that all those feet are so constituted that in each single one there is a ratio either of half again, or of double, or of equal. And so the men of whom I spoke before took account only of convenience, none of dignity.

\ciceroSection{194} For the iamb and the dactyl fall above all into verse; and therefore, just as we shun verse in prose, so these feet are to be avoided in unbroken succession — for prose is something other, and nothing is more its enemy than verses. The paean, on the other hand, is least suited to verse, and for that reason prose has welcomed it the more gladly. Ephorus, indeed, does not even understand that the spondee, which he shuns, is equal to the dactyl, which he approves; for he reckons that feet are to be measured by syllables, not by intervals — and he does the same with the trochee, which in its times and intervals is equal to the iamb, yet is faulty in prose if it is set at the close, because words fall better into longer syllables. And these same things which are in Aristotle are said about the paean by Theophrastus and Theodectes.

\ciceroSection{195} For my own part, I hold that all the feet are, as it were, mingled and blended together in prose; for we could not escape notice if we always used the same ones, since prose ought neither to be rhythmical, like a poem, nor outside rhythm, like the talk of the crowd — the one is too tightly bound, so that it appears done deliberately, the other too loosely slack, so that it seems aimless and common; by the one you would not be delighted, the other you would detest.

\ciceroSection{196} Let prose, then, as I said above, be mingled and tempered with rhythms, neither slack nor wholly rhythmical, governed by the paean above all, since the best authority so judges, but tempered also with the other rhythms which he passes over. Now it must be said which rhythms ought to be blended with which — like blending purple into a fabric — and, further, which are most suited to which kinds of speech. For the iamb is commonest in matters delivered in a lowered and humble style;

\ciceroSection{197} the paean in the loftier; and in both the dactyl. And so in prose that is varied and continuous these are to be mingled and tempered with one another. In this way the angling for delight, and the labour of squaring the prose, will least be noticed; and it will lie all the more hidden if we also draw upon the weights of words and of thoughts. For those who listen mark these two things and reckon them pleasing to themselves — words, I mean, and thoughts — and while with attentive minds they take them in with admiration, the rhythm escapes them and flies past unobserved; and yet, if it were absent, those very things would give less delight.

\ciceroSection{198} Nor indeed is the running of rhythms such — of prose, I mean, for it is far otherwise in verse — that nothing is done outside measure; for that would be a poem; rather, prose is held rhythmical when it advances all of it neither limping nor, so to speak, billowing, but evenly and steadily. And that is reckoned rhythmical in speaking which consists not wholly of rhythms, but which comes closest to rhythms; and for this reason it is harder to handle prose than verse, because in verse there is a certain fixed and definite law that must be followed, whereas in speaking nothing is laid down except that the prose be not unmeasured or cramped or slack or runaway. And so there are in it no beats of percussion, as for a flute-player, but the whole compass and contour of the prose is shut off and bounded — a thing judged by the pleasure of the ear.

\ciceroSection{199} It is usually asked, however, whether the rhythms are to be maintained throughout the whole circuit of the words, or in the first parts and in the last; for most think it enough that the sentence should fall and end rhythmically. But the truth is that this should be done supremely, not only there; for that circuit is to be set in place, not flung away. Therefore, since the ears always wait for the end and find their rest in it, that end ought not to be void of rhythm; rather, that compass of words should be borne toward this close from the very beginning, and should flow on as a whole from its head, so that, coming to the end, it may itself come to a stand.

\ciceroSection{200} And for those trained in a good discipline, who have written much and have made even whatever they say without writing resemble what is written, this will not be very hard. For the thought is first marked out in the mind, and at once the words come running up, which that same mind — than which nothing is swifter — straightway sends out, so that each may answer to its own place; and the marshalled order of these is closed off, now with one ending, now with another. And all those words, both first and middle, ought to look toward the last.

\ciceroSection{201} For at times the running of prose is more headlong, at times the advance is measured, so that from the very beginning one must see by what road one wishes to come to the end. And in rhythms no more than in the other ornaments of prose — though we do the same things the poets do — we nevertheless escape, in prose, the likeness of a poem. For in each there is both material and handling: the material is in the words, the handling in the placing of the words. And of each there are three parts: of words, the transferred, the new, the old — for of words in their proper sense I say nothing in this place; of placing, those we have named — arrangement, symmetry, rhythm.

\ciceroSection{202} But in both the poets are more frequent and freer; for they transfer words both oftener and more boldly, and use old words more readily and new ones more freely. The same happens in rhythms, in which they are compelled, as it were, to obey necessity. And yet one may understand that these arts are neither too far apart nor in no way joined. So it comes about that rhythm shows itself not in prose as it does in verse, and that what is called rhythmical in prose does not always come about by rhythm, but sometimes by the symmetry or the construction of words.

\ciceroSection{203} Thus, if the question is what the rhythm of prose is — it is all rhythm, but one kind is better and apter than another; if where its place is — in every part of the words; if from what source it springs — from the pleasure of the ear; if the method of composing it — that will be told in another place, since it belongs to practice, which was the fourth and last part of our division; if to what end it is brought into use — to give delight; if when — always; if in what place — in the whole continuous sweep of the words; if what it is that produces the pleasure — the same as in verses, whose measure the art marks out, but which the ears themselves define by a silent feeling, without art.

\ciceroSection{204} Enough has been said about the nature of the thing; practice follows, on which we must dispute more closely. Here it has been asked whether rhythm is to be maintained throughout that whole circuit of prose — which the Greeks call a \textit{[Greek: periodon]} (period), and we call now a circuit, now a compass, now a rounding-off, or a continuation, or a circumscription — or only at the beginnings, or only at the ends, or in both parts; and then, since rhythm seems to be one thing and the rhythmical another, what the difference is.

\ciceroSection{205} And further, whether it is fitting to cut the segments off in all rhythms equally, or to make some shorter and others longer, and when this should be done or why; and with which parts — with several or with single ones, unequal or equal — and when one is to use these or those; and which are most aptly placed together, and in what way; or whether there is no distinction at all in that kind; and — what bears most upon the matter — by what method prose is made rhythmical.

\ciceroSection{206} It must also be explained whence the form of the words arose, and it must be said how long the circuits should be made, and there must be discussion of their parts and, so to speak, their cuttings, and inquiry whether there is one kind and length of them, or several, and, if several, in what place or when or of what sort each ought to be used. Last of all, the usefulness of the whole kind must be set out — which indeed extends rather widely, for it is adapted not to one single thing, but to several.

\ciceroSection{207} And it is open to us, without answering point by point, so to speak of the kind as a whole that the single points too may seem to have been sufficiently answered. Setting aside, then, the other kinds, we have chosen out this one alone, which is concerned with cases and the Forum, to speak of. Now in the others — that is, in history and in what we call \textit{[Greek: epideiktikon]} (the display style) — it is approved that everything be said in the manner of Isocrates and Theopompus, with that rounding-off and compass, so that the prose may run as though enclosed in a circle, until it comes to rest upon thoughts each completed and brought to its close.

\ciceroSection{208} And so, after this — call it circumscription, or compass, or continuation, or circuit, if one may so speak — was born, no one of any standing has written prose of that kind which is framed for delight and removed from trials and the contests of the Forum without bringing nearly all his thoughts into square and rhythm. For when the hearer is one who has no fear that his trust is being assailed by the snares of well-arranged prose, he is even grateful to the orator for serving the pleasure of his ears.

\ciceroSection{209} This kind of prose, however, is neither to be taken up entire for forensic cases nor altogether to be rejected; for if you use it always, then it brings on satiety, and besides, even the untrained recognise it for what it is; moreover it takes away the feeling of the plea, it carries off the human sympathy of the hearer, it utterly destroys truth and conviction. But since it must sometimes be employed, the first thing to see is in what place, then how long it is to be kept up, and then in how many ways it is to be varied.

\ciceroSection{210} Rhythmical prose, then, is to be employed if something is to be praised more ornately — as I praised the bounty of Sicily in the second book of the prosecution, and as in the Senate I spoke of my consulship — or if a narrative is to be set out which calls for dignity rather than pathos, as in the fourth book of the prosecution I spoke of the Ceres of Henna, the Diana of Segesta, the situation of Syracuse. Often, too, in amplifying a subject, prose pours forth, by everyone's leave, rhythmical and fluent. This perhaps I have not brought to perfection, though I have very often attempted it; and that we wished and strove in mind to do so, our perorations in very many places declare. But this has force only when the man who listens is already held fast by the orator and in his grip. For then he is not bent on lying in wait and watching, but already favours him, and wishes him on, and, admiring the power of the speaking, does not search out what to blame.

\ciceroSection{211} This shape, however, is not to be maintained for long — and I do not mean in the peroration, which by its nature encloses such a structure, but in the remaining parts of a speech. For once you have made use of those passages in which I have shown that it is permissible, the whole manner of expression must be shifted over to those units which — though the Greeks call them \textit{[Greek: kommata]} and \textit{[Greek: kōla]} — we, I cannot think why, less than correctly call phrases and clauses. For there can be no familiar names for things that are unfamiliar; but since we are accustomed, whether for the sake of charm or for want of a term, to transfer words from one use to another, this happens in all the arts: when a thing must be named which, owing to the very ignorance of the things themselves, has had no name before, necessity compels us either to coin a new word or to borrow one from something similar.

\ciceroSection{212} In what way it is fitting to speak in phrases and in clauses we shall see presently; for now I must speak of how many ways periods and their closings may be varied. Rhythm, on the whole, flows now more quickly, by the shortness of its feet, now more slowly, by their length. The pace of the running is what contentious passages demand, slowness what the setting out of facts requires. The circuit comes to rest in several measures, of which there is one that Asia in particular followed — the one called the double trochee, when the two final feet are trochees, that is, each made of a long and a short. (For it must be explained that the same feet are named by different writers under different terms.)

\ciceroSection{213} The double trochee is not, indeed, of itself faulty in the close; but in the rhythm of prose nothing is so faulty as for it always to be the same. And that very foot falls of itself admirably — for which reason satiety is all the more to be dreaded. When I myself was standing by, C. Carbo, son of Gaius, tribune of the plebs, said this in an assembly, in these words: \textit{O Marcus Drusus, I appeal to your father} — these two, indeed, in two phrases, each of two feet; then by clauses: \textit{You used to say that the commonwealth is a sacred thing} — these likewise are clauses, of three feet each; then a circuit:

\ciceroSection{214} \textit{that whoever had violated her had paid the penalty in full to all} — a double trochee; for it makes no difference to the matter whether that final syllable is long or short; then: \textit{The father's saying the son's rashness has proved true} — by this double trochee so great an outcry was roused from the assembly that it was a marvel. I ask: is it not the rhythm that achieved this? Change the order of the words, make it run thus: \textit{has proved true the son's rashness}, and at once there will be nothing — even though "rashness" is made of three shorts and a long, the very measure which Aristotle approves as best, in which I disagree with him. "But the words are the same, the sense the same."

\ciceroSection{215} For the mind that is enough; for the ears it is not enough. But this should not be done too often: for first the rhythm is recognized, then it cloys, and afterward, once the trick of it is understood, it is held in contempt. Yet there are several closes that fall rhythmically and agreeably. For the cretic, which is made of a long, a short, and a long, and its equivalent the paean, which is equal in span though longer by a syllable, is thought to be very conveniently bound into loose prose, since it is twofold. For it is made either of a long and three shorts, a measure which thrives at the opening but lies flat at the close, or of as many shorts and a long, the one into which the ancients judge it falls best; I do not flatly reject this, but I prefer others.

\ciceroSection{216} Not even the spondee is to be utterly repudiated, although, since it is made of two longs, it seems duller and slower; still, it has a certain steady step, not without dignity — and in phrases far more so, and in clauses; for it makes up for its scanty count of feet by its own gravity and slowness. But when I name these feet in the close, I am not speaking of a single final foot: I add — and this is the least — the one immediately before, and often a third as well.

\ciceroSection{217} Not even the iamb, which is made of a short and a long, or the trochee, which is equal to the foot of three shorts but equal in span, not in syllables, or even the dactyl, which is made of a long and two shorts — if it stands next to the last foot — arrives at the close with too little smoothness, provided the final foot is a trochee or a spondee; for it never matters which of the two stands in the final foot. But these same three feet make a poor conclusion if any of them is placed at the very end, except when the last is a dactyl in place of a cretic; for it makes no difference whether the dactyl is last or the cretic, since whether the final syllable is short or long does not matter even in verse.

\ciceroSection{218} For which reason the man who said the paean was the more fitting on the ground that its last syllable was long saw too little, since it makes no difference how long the last syllable is. And further, the paean, because it has more than three syllables, is reckoned by some a rhythm, not a foot. It is indeed — as all the ancients agree, Aristotle, Theophrastus, Theodectes, Ephorus — the one most fitted to a speech, whether at its opening or in its middle; those writers think it fit even at the falling close, where the cretic seems to me the more fitting. The dochmiac, however, made of five syllables — a short, two longs, a short, a long, as in this case: \textit{You hold your friends} — is fit in any position, provided it be set down only once: repeated or run on, it makes the rhythm too plain and too conspicuous.

\ciceroSection{219} If, then, we make use of these changes, so many and so varied, neither will it be detected outright that anything is being done by us on purpose, and satiety will be forestalled. And since rhythmic prose is made not by rhythm alone but also by the arrangement and by that kind of symmetry which has been spoken of before — by arrangement it can be understood, when the words are so built that the rhythm seems not sought after but to have followed of its own accord, as in Crassus: \textit{For where lust holds sway, innocence has but slight protection}; here the order of the words produces the rhythm without any open effort on the orator's part — and so, if those ancients, I mean Herodotus and Thucydides and that whole age, said anything aptly and rhythmically, it fell out so not by rhythm sought after but by the placing of the words.

\ciceroSection{220} There are, moreover, certain shapes of speech in which the symmetry is such that rhythm follows of necessity. For when like is balanced against like, or contrary set against contrary, or words answering to words are paired with similar endings, whatever is so concluded for the most part falls out rhythmically — a kind of which we have spoken above, with examples — so that this resource too may afford the means of not always ending in the same way. And yet these things are not so cramped and bound that we cannot loosen them when we wish. There is a great difference between prose that is rhythmic — that is, like the rhythms — and prose that consists plainly of rhythms; if the latter happens, it is an intolerable fault; if the former does not happen, the prose is scattered and uncultivated and runs to waste.

\ciceroSection{221} But since it is not only not frequently but even rarely, in real cases or in forensic ones, that one ought to speak in rounded and rhythmic periods, it seems to follow that we should consider what those things are which I called above phrases and clauses. For in real cases these hold the greatest part of a speech. The circuit and the full period consists, as a rule, of four parts, which we call clauses, in such a way that it fills the ear and is neither shorter than is sufficient nor longer. And yet both sometimes happen — or rather, often — that one must either stop sooner or proceed further, lest brevity seem to have cheated the ear or length to have dulled it. But I keep an account of the mean; for I am not speaking of verse, and prose is somewhat freer. The full period, then, consists of about four lines, as it were on the scale of hexameter verses.

\ciceroSection{222} In these single lines, then, there appear, as it were, the knots of the continuation, which we join together in the circuit. But if we wish to speak by clauses, we come to rest, and, when there is need, we readily and often disengage ourselves from that invidious running. But nothing ought to be so rhythmic as this, which least appears and avails the most. Of this kind is that passage of Crassus: \textit{Let them dismiss their advocates; let them come forward themselves} — had he not set \textit{let them come forward themselves} at an interval, he would surely have felt that he had poured out a senarius; on the whole it would fall out better as \textit{themselves come forward}; but it is the kind of thing I am now discussing —

\ciceroSection{223} \textit{Why do they assail us with secret designs? Why do they muster forces against us out of our own deserters?} The first two are those which the Greeks call \textit{[Greek: kommata]}, we phrases; then the third is what they call a \textit{[Greek: kōlon]}, we a clause; there follows one not long — for the period is completed out of two lines, that is, clauses, and falls into spondees; and Crassus, indeed, for the most part spoke in this way, and this kind of speaking I myself most approve. But the things that are brought out in phrases or in clauses ought to fall most aptly of all, as in this passage of mine: \textit{Did you lack a house? But you had one. Had you money to spare? But you were in want;}

\ciceroSection{224} these were spoken in four phrases; but the two that follow are by clauses: \textit{Out of your senses, you ran headlong into the columns; out of your mind, you raged against what was not yours.} Then everything is held up, as if on a kind of base, by a longer period: \textit{You valued your house, sunk and blind and prostrate, at more than yourself and more than your fortunes.} It is closed with a double trochee. But the foregoing closes with spondees. For in those passages, which one ought to use as little daggers, brevity itself will make the feet freer; for one must often use single feet, for the most part two, and to each may be added part of a foot, but generally not more than three.

\ciceroSection{225} A speech handled in phrases and in clauses avails the most in real cases, and most of all in those passages where you either bring a charge or refute one, as I did in the second speech for Cornelius: \textit{O the cunning men! O the contrivance! O the genius to be feared!} — by clauses thus far; then by cuts: \textit{We said it}; again by clauses: \textit{We are willing to produce witnesses.} The final period follows, but out of two clauses, than which it cannot be shorter: \textit{Which of us, I ask, did it deceive that you would act so?}

\ciceroSection{226} Nor is there any kind of speaking either better or stronger than to strike with two or three words, sometimes with single ones, a little elsewhere with more, among which, with varied closes, a rhythmic period interposes itself rarely; and Hegesias, fleeing it perversely, while he too wishes to imitate Lysias, almost a second Demosthenes, dances as he chops up his little particles. And that man indeed offends no less in his thoughts than in his words, so that whoever has come to know him need not look about for someone to call inept. But I have set down those passages of Crassus and my own, that anyone who wished might judge by the ears themselves what is rhythmic even in the smallest particles of a speech. And since we have said more about rhythmic prose than anyone before us, we shall now speak of the usefulness of this kind.

\ciceroSection{227} For there is nothing else, Brutus — as you, least of all men, are unaware — in speaking beautifully and like an orator than to speak with the best thoughts and the choicest words. And no thought brings the orator any fruit unless it is aptly and fully set out, nor does the brilliance of words appear unless they are carefully placed, and rhythm sets both of these in their light; yet rhythm — for this must often be attested — is not only not poetically bound, but even flees that bondage and is most unlike it of all things; not that the rhythms are not the same, not only of orators and of poets but of all who speak, indeed of all sounds whatever that we can measure by the ear, but the order of the feet makes what is uttered seem like prose or like poetry.

\ciceroSection{228} This, then — whether it please you to call it arrangement, or finish, or rhythm — must necessarily be employed, if you would speak with ornament; and not only, as Aristotle and Theophrastus say, lest the speech be carried on without limit like a river, which ought to come to rest not by the breath of the speaker or by the punctuation of the copyist but as compelled by rhythm, but also because what is apt has far greater force than what is loose. For just as we see that athletes — and gladiators not much otherwise — do nothing either in dodging warily or in striking vehemently in which this motion does not have a certain grace of the wrestling-school, so that whatever is done usefully in these matters for the fight is at the same time comely to behold, so the orator neither lands a heavy blow unless his thrust was apt, nor turns aside an onset closely enough unless even in giving way he understands what is fitting.

\ciceroSection{229} And so, such as is the motion of those whom the Greeks call ungainly \textit{[Greek: apalaistrous]}, such seems to me the speech of those who do not close their thoughts with rhythms; and so far is it from true — what those who, whether through want of teachers or slowness of talent or flight from labour, have not attained this, are wont to say, that speech is enervated by the arrangement of words — that on no other terms can there be in it any force or onset at all. But the thing demands great practice, lest we do anything like those who, having pursued this kind, did not master it; lest we transpose words openly, in order that the speech may fall or roll the better; which

\ciceroSection{230} L. Caelius Antipater, in the preface to his Punic War, declares he will not do unless under necessity. O the simple man, who hides nothing from us — the wise man, who thinks one must be a slave to necessity! But he is altogether unschooled; for our part, in writing and in speaking, the excuse of necessity does not win our approval; for nothing is necessary, and if anything were, it still was not necessary to confess it. And he indeed, who begs this indulgence from Laelius, to whom he wrote and to whom he excuses himself, both uses that transposition of words and yet none the more aptly fills out and concludes his thoughts. Among others, and especially among the Asiatics who are slaves to rhythm, you will find certain empty words crammed in, like fillers for the rhythms. There are also those who, through that fault which flowed chiefly from Hegesias, by breaking and chopping up their rhythms fall into a kind of low style very like that of little verses.

\ciceroSection{231} There is a third kind, in which were those famous brothers, the leaders of the Asiatic rhetors, Hierocles and Menecles, by no means, in my judgment, to be despised. For although they are far from the form of truth and from the standard of the Attic writers, they yet make up for this fault by their facility or their fullness. But among them there was no variety, since almost everything was concluded in one way. Whoever has fled these faults, so that he neither transposes a word in such a way that it is understood to have been done on purpose, nor stuffs in words to fill, as it were, the chinks, nor, pursuing minute rhythms, chops up and unstrings his thoughts, nor without any change keeps always turning in the same kind of rhythms — he will have avoided nearly all the faults. For of the merits we have said much, to which those faults are plainly contrary.

\ciceroSection{232} How much it matters to speak aptly may be tested if you dissolve the well-built placing of a finished orator by transposing the words; for the whole thing is spoiled, as both this passage of mine in the speech for Cornelius and all that follows: \textit{Nor do riches move me, in which many a slave-dealer and trafficker has surpassed all the Africani and Laelii}: change it a little, so that it runs \textit{many traffickers and slave-dealers have surpassed}, and the whole thing will have perished; and what follows: \textit{Nor the raiment or the chased gold and silver in which many a eunuch out of Syria and Egypt has outdone our old Marcelli and Maximi}; transpose the words thus, so that it runs \textit{have outdone the eunuchs out of Syria and Egypt}; add the third: \textit{Nor indeed those ornaments of country villas, by which I see that L. Paullus and L. Mummius, who crammed the city and all Italy with such things, could very easily have been surpassed by some man of Delos or Syria}; make it thus: \textit{could have been surpassed by some Syrian or man of Delos};

\ciceroSection{233} do you see how, the order of the words being changed a little, while the sense stands in the very same words, everything falls away to nothing, once the apt has been dissolved? Or, if you take from some ill-arranged writer a scattered thought and, by changing the order of the words a little, bring it into square, that will be made apt which before was loose and flowing away. Come, take that saying of Gracchus before the censors: \textit{It cannot be but that it belongs to the same man to disapprove the upright who approves the base}; how much more aptly, had he said it thus: \textit{But that it belongs to the same man, who approves the base, to disapprove the upright!}

\ciceroSection{234} No one has ever been unwilling to speak in this manner, nor has anyone been able to do so without speaking thus; but those who spoke otherwise were not able to attain this. So they became Attics on the sudden — as if Demosthenes had been a man of Tralles! whose famous thunderbolts would not flash so, were they not hurled, twisted home, by rhythms. But if loose prose delights anyone the more, let him by all means pursue it — only on this condition: that, as if one had dissolved the shield of Phidias, he will have destroyed the whole effect of the placing, not the loveliness of the individual works; just as in Thucydides I only miss the rounded shape of the speech — the ornaments are there to be seen.

\ciceroSection{235} But those men, when they dissolve a speech in which there is neither matter nor a single word that is not low, seem to me to be dissolving not a shield, but — as the proverb has it (though it is more humbly put, it is yet apt) — a besom, so to speak. And that they may seem plainly to have despised this kind which I praise, let them write something either in the manner of Isocrates or in that which Aeschines or Demosthenes uses; then I shall judge that they shrank from this kind not out of despair but recoiled from it by judgment; or I shall myself find someone willing to use it on the same terms, so as to speak or write in whichever language you please, in the very kind they prefer; for it is easier to dissolve what is apt than to connect what is scattered.

\ciceroSection{236} The matter, however, stands thus, that I may say most briefly what I think: to speak with arrangement and aptly, but without thoughts, is madness; to speak with weight of thought, but without order and measure in the words, is to be a babbler — yet a babbler of such a kind that those who practise it can be reckoned not foolish men, even, for the most part, prudent ones; and whoever is content with that, let him use it. But the eloquent man, who ought to move not only approvals but admiration, outcries, applause, if it be permitted, must so excel in all things that it would be a disgrace to him for anything to be either expected or heard with more pleasure.

\ciceroSection{237} You have my judgment concerning the orator, Brutus; either follow it, if you approve it, or hold to your own, if your own is something other. In this I shall neither contend with you, nor ever affirm that this judgment of mine, which I have so strongly maintained in this book, is truer than yours. For not only may one thing seem so to me and another to you, but to my very own self one thing may seem so at one time and another at another. And not only in this matter, which looks to the assent of the crowd and to the pleasure of the ear — the two lightest things in judging — but not even in the greatest matters have I yet found anything firmer to hold to, or by which to direct my judgment, than whatever seemed to me most like the truth, since the truth itself nevertheless lies hidden in concealment.

\ciceroSection{238} But I should wish, if those things which have been argued win your approval too little, that you either think a greater work was undertaken than could be accomplished, or that, while I wished to comply with your asking, out of shame at refusing I took upon myself the impudence of writing.
